% Protocol V3 manuscript body. Numeric claims are generated from locked test summaries.

\section{Background}
\label{sec:intro}
Vision--language pretraining has made free-text conditioning practical for dense
prediction \cite{radford2021clip,clipseg,biomedclip}. In medical segmentation,
language can identify a target, describe morphology, or provide spatial context
\cite{lvit,causalclipseg,medclipseg}. This additional modality also creates two
reliability requirements that do not exist in image-only segmentation: training
transforms must preserve the joint meaning of the image, mask, and prompt, and
the resulting prediction should depend on the case-matched prompt when that
prompt is claimed to provide useful guidance.

Conventional geometric augmentation applies one transform to the image and mask
but leaves text unchanged \cite{shorten2019survey}. If a prompt locates a lesion
in the image frame, a horizontal flip changes \textit{left} to \textit{right},
a vertical flip changes \textit{upper} to \textit{lower}, and a right-angle
rotation changes both axes. The transformed mask remains geometrically valid,
yet the conditioning text may become false. This is an avoidable cross-modal
supervision conflict. Separately, a geometrically valid prompt can still be
uninformative or weakly grounded; high segmentation accuracy alone cannot show
that the text branch is being used. The relevant augmentation example is
image-frame location, rather than
anatomical laterality: mirroring a breast image does not change the patient's
anatomical left or right.

We therefore treat the shared training policy as the main empirical object and
prompt reliability as a control on how its accuracy is obtained. \methodname{}
combines biomedical text features, adaptive convolution, and synchronized
image--mask--text augmentation, while matched and fixed-checkpoint controls
prevent the aggregate result from being misattributed to rewriting or to
clinically meaningful language understanding. Our contributions are:
\begin{enumerate}
  \item A complete retrospective 90-cell ordered component ablation that
  separates the base pathway, frozen pooled BioMedCLIP embedding adapter, the
  effective two-projector ATConv adapter, the shared training plan, and prompt
  rewriting.
  \item A same-data-and-evaluation-protocol comparison on five public datasets,
  with three training seeds, restored native-grid mask scoring, and three
  text-guided and six image-only architectures.
  \item A protocol and prompt-reliability audit comprising an exact-plan
  no-rewrite control, correct/shuffled/fixed/empty inference controls,
  a post hoc presence/class-stratified prompt derangement, constant-text and
  Base-plus-plan follow-ups, prompt-provenance boundaries, and a runtime check
  of which registered ATConv modules enter the forward graph.
\end{enumerate}

\subsection{Related work}
\subsubsection{Text-guided segmentation}
CLIP learns a transferable joint image--text space \cite{radford2021clip}, and
CLIPSeg adapts it to dense prediction through a lightweight conditional decoder
\cite{clipseg}. ViLT shows that joint vision--language interaction can be built
without a region detector or convolutional visual front end \cite{vilt}, while
BioMedCLIP introduces biomedical image--text pretraining at scale
\cite{biomedclip}. LViT integrates language with a transformer--UNet pathway
\cite{lvit}; CausalCLIPSeg uses causal intervention for referring medical
segmentation \cite{causalclipseg}; and MedCLIPSeg uses probabilistic cross-modal
adaptation \cite{medclipseg}. TextMatch studies multimodal consistency in
semi-supervised segmentation \cite{textmatch}, whereas ARSeg targets missing or
imbalanced textual attributes \cite{arseg}. These models motivate a controlled
test of whether the training prompt still describes the transformed image.

Recent work also separates prompt content from fusion design. DescriptorMedSAM
evaluates a structured progression from target name to name--shape--location
descriptors and finds that shape and location provide complementary guidance
\cite{descriptormedsam}. Spatial-Aware Symmetric Alignment encodes directional
constraints as region-level guidance while aligning descriptive and diagnostic
expressions \cite{ssa}. NTP-MRISeg instead casts image, text, and mask tokens as
one autoregressive sequence \cite{ntpmriseg}. These architecture- and
descriptor-level advances are complementary to \methodname{}, which asks both
whether a supplied spatial descriptor remains valid after training-time geometry
and whether the trained prediction remains case-specifically text dependent.

\subsubsection{Image-only reference models}
U-Net \cite{unet}, U-Net++ \cite{unetpp}, nnU-Net \cite{nnunet}, Rolling-Unet \cite{rollingunet},
U-Mamba \cite{umamba}, and U-KAN \cite{ukan} demonstrate that a strong visual
pathway can perform well without text. A multimodal method should therefore be
compared against both text-guided baselines and competitive image-only models.
RollingUNet, U-Net, U-Net++, nnU-Net, U-Mamba, and U-KAN are complete for all
datasets and all three seeds under Protocol V3. This breadth is important
because the strongest image-only reference differs across datasets.

\subsubsection{Consistency under paired augmentation}
Data augmentation generally assumes that the target remains valid after a
transformation \cite{shorten2019survey}. For image--text learning, this assumption
extends to language. Yi et al. update predefined caption keywords when
horizontal flipping image--caption pairs and combine this with
semantics-preserving appearance augmentation \cite{augmentpairs}. In multimodal
medical segmentation, RBTCA instead grounds textual cues in image regions before
augmenting a modality-aware representation \cite{rbtca}. Chai et al. identify the same flip/rotation conflict and
avoid it by fusing text into a pseudo-image before augmentation
\cite{chai2025earlyfusion}. Semi-MedRef synchronizes positional spans during
patch mixing and perturbs position phrases in a semi-supervised teacher--student
setting \cite{semimedref}. \methodname{} therefore does not claim priority for the
general alignment problem, paired caption augmentation, or exchanging individual
left/right tokens. Its specific contribution is the cross-dataset evidence chain:
an explicit lexical map for a finite image-coordinate vocabulary, deterministic
exact-plan controls, a complete ordered ablation, and a privileged-prompt
reliability audit.

\section{Methods}
\label{sec:protocol}
\subsection{Dataset and prompt provenance}
Table~\ref{tab:bspc_dataset_audit} reports the locked manifests. The image and
mask sources are BUSI \cite{busi}, CVC-ClinicDB \cite{cvc}, BUS-BRA
\cite{busbra}, BRISC \cite{brisc}, and QaTa-COV19 \cite{qatacov19}. Their paired
prompts are distributed by the MedCLIPSeg repository \cite{medclipseg}. For
lesion cases, MedCLIPSeg extracts class, coarse location, component count, shape,
and tone from the ground-truth mask or label and fills language-model-generated
templates. Thus these are target-derived synthetic referring captions, not
prospective clinical reports. This benchmark construction is shared by the
image+text methods in Table~\ref{tab:bspc_v3_models}, but it gives them
target-derived information unavailable to image-only methods and limits both
clinical interpretation and cross-modality comparison.

\def\BmcTableSelection{1}
% Tables for the BMC Medical Imaging manuscript.
% Ordered by the evidence hierarchy: dataset audit, ordered configurations,
% semantic controls, contextual model comparison, and exploratory diagnostics.

\ifnum\BmcTableSelection=1
\begin{table}[htbp]
\centering
\caption{Image-level dataset partitions and realized prompt-change frequencies.}
\label{tab:bspc_dataset_audit}
\scriptsize
\setlength{\tabcolsep}{2.0pt}
\begin{tabular}{@{}lrrrlrrr@{}}
\toprule
Dataset & Train & Val. & Test & Prompt & Unique & Spatial (\%) & Change (\%) \\
\midrule
BUSI     & 624  & 78  & 78   & Target-derived & 198 & 82.85 & $20.27\pm0.10$ \\
ClinicDB & 490  & 61  & 61   & Target-derived & 202 & 100.00 & $28.48\pm0.21$ \\
BUS-BRA  & 1311 & 282 & 282  & Target-derived & 75  & 100.00 & $12.44\pm0.12$ \\
BRISC    & 3881 & 974 & 892  & Target-derived & 359 & 78.23 & $5.75\pm0.04$ \\
COVID-19 & 5716 & 1429& 2113 & Target-derived & 269 & 100.00 & $45.75\pm0.06$ \\
\bottomrule
\end{tabular}
\vspace{2pt}
\begin{minipage}{0.98\linewidth}
\footnotesize\raggedright
\textit{Notes:} Unique and spatial coverage use training prompts. Realized
change enumerates all case--epoch events under the fixed 100-epoch shared
augmentation policy and reports mean$\pm$sample standard deviation across seeds
123, 456, and 789. Public prompts are target-derived synthetic captions, not
clinical text.\par
\end{minipage}
\end{table}
\fi

\ifnum\BmcTableSelection=3
\begin{table}[htbp]
\centering
\caption{Inference-time prompt controls for trained \methodname{} models on
five public datasets.}
\label{tab:bspc_semantic_controls}
\scriptsize
\setlength{\tabcolsep}{4.5pt}
\begin{tabular}{@{}lrrrrr@{}}
\toprule
Dataset & True Dice & Shuffled Dice & Fixed Dice & Empty Dice & $\Delta$ pp [95\% CI] \\
\midrule
BUSI     & $85.78\pm0.45$ & $46.73\pm1.41$ & $16.69\pm0.02$ & $13.22\pm5.35$ & $+39.05$ [$29.85,48.49$] \\
ClinicDB & $92.07\pm0.71$ & $90.20\pm0.94$ & $88.31\pm1.31$ & $78.71\pm7.93$ & $+1.88$ [$0.07,4.22$] \\
BUS-BRA  & $90.24\pm0.14$ & $89.47\pm0.23$ & $88.87\pm0.28$ & $68.22\pm8.44$ & $+0.77$ [$0.45,1.15$] \\
BRISC    & $87.18\pm0.14$ & $70.89\pm2.12$ & $75.31\pm1.97$ & $58.89\pm5.30$ & $+16.29$ [$14.26,18.43$] \\
COVID-19 & $84.08\pm0.13$ & $68.80\pm0.82$ & $68.54\pm0.16$ & $55.81\pm3.82$ & $+15.29$ [$14.49,16.10$] \\
\bottomrule
\end{tabular}
\vspace{2pt}
\begin{minipage}{0.98\linewidth}
\footnotesize\raggedright
\textit{Notes:} No model is retrained; only the test prompt changes. Dice is
mean$\pm$sample standard deviation over three seeds. The last column is true
minus shuffled Dice with a paired-case bootstrap 95\% interval after
three-seed averaging. Holm-adjusted Wilcoxon tests and the bootstrap intervals
are positive on all five datasets.\par
\end{minipage}
\end{table}
\fi

\ifnum\BmcTableSelection=4
% Auto-generated by generate_medequiseg_factorial_manuscript_assets.py.
% Do not edit numerical cells manually.
\begin{table}[htbp]
\centering
\caption{Five-dataset macro-average accuracy across five retained ordered configurations.}
\label{tab:medequiseg_factorial}
\scriptsize
\setlength{\tabcolsep}{2.4pt}
\begin{tabular}{@{}lccccrr@{}}
\toprule
Configuration & Text enc. & ATConv & Aug. & Rewrite & Dice & IoU \\
\midrule
Base model & -- & -- & -- & -- & $86.41\pm0.35$ & $79.22\pm0.34$ \\
$+$ BioMedCLIP & \checkmark & -- & -- & -- & $86.67\pm0.08$ & $79.52\pm0.06$ \\
$+$ ATConv & -- & \checkmark & -- & -- & $86.63\pm0.32$ & $79.42\pm0.31$ \\
$+$ BioMedCLIP $+$ ATConv & \checkmark & \checkmark & -- & -- & \underline{$87.03\pm0.15$} & \underline{$79.85\pm0.23$} \\
\textbf{\methodname{} (Ours)} & \checkmark & \checkmark & \checkmark & \checkmark & \textbf{$87.87\pm0.17$} & \textbf{$80.97\pm0.13$} \\
\bottomrule
\end{tabular}
\vspace{2pt}
\begin{minipage}{0.98\linewidth}
\footnotesize\raggedright
\textit{Notes:} Five retained configurations, five datasets, and three seeds produce 75
training runs. Values are five-dataset macro-average percentages
(mean$\pm$sample standard deviation). Text enc. denotes pooled BioMedCLIP
conditioning; ATConv denotes the two projection-head operators active in the evaluated forward pass;
Aug. denotes the common shared augmentation policy. Shared augmentation and
prompt rewriting appear together only in the final configuration, so these are conditional configuration contrasts,
not independent component effects or complete interactions. Bold and underline
mark the best and second-best configurations. The complete per-dataset matrix
is in Supplementary Table S9.\par
\end{minipage}
\end{table}

\fi

\ifnum\BmcTableSelection=5
\begin{table}[htbp]
\centering
\caption{Per-dataset and equal-dataset macro Dice on five public benchmarks.}
\label{tab:bspc_v3_models}
\scriptsize
\setlength{\tabcolsep}{2.0pt}
\begin{tabular}{@{}lrrrrrr@{}}
\toprule
Method & BUSI & ClinicDB & BUS-BRA & BRISC & COVID-19 & Macro avg. \\
\midrule
\multicolumn{7}{l}{\textit{Image+text methods}} \\
MedCLIPSeg
 & $\mathbf{86.13}$ & $\underline{91.96}$ & $88.47$
 & $81.36$ & $\underline{82.62}$ & $86.11$ \\
CausalCLIPSeg
 & $84.62$ & $91.19$ & $\underline{88.87}$
 & $\underline{86.11}$ & $81.85$ & $\underline{86.53}$ \\
LViT
 & $81.81$ & $91.03$ & $87.66$
 & $85.24$ & $78.50$ & $84.85$ \\
\textbf{\methodname{} (Ours)}
 & $\underline{85.78}$ & $\mathbf{92.07}$ & $\mathbf{90.24}$
 & $\mathbf{87.18}$ & $\mathbf{84.08}$ & $\mathbf{87.87}$ \\
\midrule
\multicolumn{7}{l}{\textit{Image-only contextual references}} \\
RollingUNet
 & $72.36$ & $87.31$ & $88.00$
 & $85.03$ & $79.32$ & $82.40$ \\
U-Net
 & $61.31$ & $85.61$ & $87.35$
 & $85.00$ & $78.81$ & $79.62$ \\
U-Net++
 & $76.21$ & $89.77$ & $88.59$
 & $\underline{85.94}$ & $\mathbf{80.00}$ & $84.10$ \\
nnU-Net
 & $74.99$ & $\mathbf{93.62}$ & $\underline{89.80}$
 & $\mathbf{87.30}$ & $79.07$ & $\mathbf{84.96}$ \\
U-Mamba
 & $\underline{76.63}$ & $\underline{92.21}$ & $\mathbf{90.34}$
 & $85.37$ & $79.08$ & $\underline{84.73}$ \\
U-KAN
 & $\mathbf{77.10}$ & $91.40$ & $88.56$
 & $85.45$ & $\underline{79.38}$ & $84.38$ \\
\bottomrule
\end{tabular}
\vspace{2pt}
\begin{minipage}{0.98\linewidth}
\footnotesize\raggedright
\textit{Notes:} Dataset values are three-seed mean Dice percentages; the macro
average is the arithmetic mean of the five dataset means. Bold and underline rank
methods only within the image+text
or image-only panel. Cross-panel differences are descriptive because the
public image+text methods receive target-derived synthetic captions unavailable
to image-only methods. Complete mean$\pm$sample standard deviation Dice and IoU
results are in Supplementary Table S14. The five \methodname{} entries use the
models from the ordered configuration analysis.\par
\end{minipage}
\end{table}
\fi

\ifnum\BmcTableSelection=7
\begin{table}[t]
\centering
\caption{Equal-dataset macro segmentation accuracy on the five public
benchmarks. Values are the arithmetic mean of the three-seed dataset means in
Table~\ref{tab:bspc_v3_models}. Bold and underline rank
models within modality. Cross-modality differences are descriptive because the
image+text models receive target-derived synthetic captions unavailable to the
image-only models.}
\label{tab:public5_macro_accuracy}
\small
\setlength{\tabcolsep}{6pt}
\begin{tabular}{lcc}
\toprule
Method & Dice (\%) & IoU (\%) \\
\midrule
\multicolumn{3}{l}{\textit{Image+text methods}} \\
\textbf{\methodname{} (Ours)} & $\mathbf{87.87}$ & $\mathbf{80.97}$ \\
CausalCLIPSeg Baseline & $\underline{86.53}$ & $\underline{79.32}$ \\
MedCLIPSeg & $86.11$ & $78.27$ \\
LViT & $84.85$ & $77.37$ \\
\midrule
\multicolumn{3}{l}{\textit{Image-only contextual references}} \\
nnU-Net & $\mathbf{84.96}$ & $\mathbf{78.23}$ \\
U-Mamba & $\underline{84.73}$ & $\underline{77.94}$ \\
U-KAN & $84.38$ & $77.02$ \\
U-Net++ & $84.10$ & $76.65$ \\
RollingUNet & $82.40$ & $74.62$ \\
U-Net & $79.62$ & $71.20$ \\
\bottomrule
\end{tabular}
\end{table}
\fi

\ifnum\BmcTableSelection=6
\begin{table}[htbp]
\centering
\caption{Exploratory three-seed continuous-location recomputation analysis.
\locrecompute{}
recomputes coarse position phrases from the transformed training mask when a
mild affine, translation, or zoom is sampled. Values are mean$\pm$sample
standard deviation. Delta and the reported bootstrap interval are
\locrecompute{} minus its reference model in percentage points. Both runs use
the same transformation distribution but independently initialized augmentation
sequences, so the realized per-case transformations are not paired. These
rows are not pooled with the main ordered configuration analysis. $p_W$ is the exploratory, unadjusted two-sided
paired Wilcoxon $p$-value.}
\label{tab:continuous_location_diagnostic}
\small
\setlength{\tabcolsep}{4.2pt}
\begin{tabular}{@{}llrrrr@{}}
\toprule
Dataset & Metric & Reference model & Location Recompute & $\Delta$ pp [95\% CI] & $p_W$ \\
\midrule
BUSI & Dice & $85.74\pm0.80$ & $84.62\pm0.20$ & $-1.12$ [$-2.62,0.07$] & .004 \\
     & IoU  & $79.46\pm0.83$ & $78.06\pm0.64$ & $-1.40$ [$-2.70,-0.30$] & .002 \\
ClinicDB & Dice & $93.18\pm0.54$ & $93.06\pm0.72$ & $-0.13$ [$-1.25,1.27$] & .138 \\
         & IoU  & $87.77\pm0.82$ & $87.54\pm1.07$ & $-0.23$ [$-1.94,1.78$] & .117 \\
\bottomrule
\end{tabular}
\end{table}
\fi


\subsection{Prompt-evidence scope and real-text validation}
The public MedCLIPSeg mask/label-derived captions are restricted to mechanism
analysis. They do not constitute real-text validation. Deployment-oriented
evidence, which is not available in the present study, would require blinded
clinician-authored or clinician-verified prompts generated from images without
access to masks. Such a study must freeze the vocabulary and annotation
instructions before opening the untouched test set and must compare correct,
shuffled, fixed, and empty prompts. We do not use machine-generated captions as
a substitute for clinical text evidence.

Spatial coverage counts prompts containing any member of the finite dictionary,
including center, middle, top, bottom, upper, lower, left, and right compounds.
This count alone overstates intervention frequency because central phrases can be
invariant and geometric operations are sampled probabilistically. We therefore
replay the exact keyed random generator for every training case, epoch, and seed
and report the realized proportion of events whose final prompt string differs
from the original. Across 100 epochs and seeds 123, 456, and 789, the rates
range from $5.75\pm0.04\%$ for BRISC to $45.75\pm0.06\%$ for COVID-19.
Complete event counts and operation rates are reported in Supplementary
Table~S7.

\subsection{Protocol V3}
The protocol identifier is \texttt{MEDSEG\_TEXT\_V3\_20260710}. Manifests are
fixed across seeds. Seeds 123, 456, and 789 change parameter initialization,
minibatch order, and augmentation sampling; they do not resplit the data.
Validation selects the best checkpoint, and the test split is accessed only by
the final prediction and evaluation scripts. Predictions are restored to the
native ground-truth resolution before scoring.

The protocol hash locks the dataset registry, manifests, split roles, mask
rules, seeds, threshold, and evaluation specification, but it does not encode
the full training policy or evaluator source. In particular, the original V3 LCAugv2 router restricted ClinicDB and
COVID-19 to horizontal flips and excluded affine, translation, and zoom. The
shared-plan rows in the 90-cell matrix instead use the separately dated R11
follow-up policy: all five public datasets receive the same multi-geometry,
appearance, mild affine, translation, and zoom sampler. That policy was frozen
before the paired R11/no-rewrite matrix, but it is an amendment on the fixed V3
data and evaluation protocol rather than part of the original protocol lock.
Accordingly, we use the V3 protocol hash only for data/evaluation identity and
retain a separate run-configuration hash, code hash, manifest hash, cache hash,
checkpoint hash, and original command log for each training task. We do not
describe the shared-plan matrix as prospectively confirmatory under the original
July 10 augmentation specification (Supplementary Table S11).

BRISC contains duplicated images with conflicting masks. Protocol V3 drops
conflicting duplicate groups while preserving exact holdout membership; this
sanitation rule is part of the dataset registry and protocol hash. The public
MedCLIPSeg manifests inherit the repository's
Train/Val/Test folders and do not expose a populated patient- or sequence-ID
field; consequently, we describe these as fixed image-level holdouts rather than
claiming patient-level separation. A post hoc filename audit can recover subject
IDs for part of COVID-19: no recovered subject crosses into the fixed test set,
but 434 subjects occur in both its original training and validation folders. We
therefore performed a separately hashed sensitivity rerun that assigns every
recoverable subject to only one training or validation split while preserving the
fixed test set. Rows without recoverable IDs remain case-level groups. The
grouped-validation training results are reported in Supplementary Table S10, and
the reported-checkpoint test-set cluster sensitivity additionally addresses
uncertainty from repeated test images.

We report Dice and IoU in the main table and additionally export precision,
recall, normalized surface Dice (NSD) at two pixels, HD95, and ASSD. A both-empty
case receives perfect overlap and zero distance. A one-sided empty case receives
zero overlap and infinite distance, and the number of such boundary failures is
reported with the distance metrics. Supplementary Tables S1--S5 report NSD,
finite-case HD95 and ASSD, and the one-empty failure rate for every method and
dataset. Finite distance means are interpreted only together with the failure
rate.

\subsection{MedEquiSeg framework}
\label{sec:method}
\subsubsection{Framework overview and reliability scope}
\methodname{} comprises a visual segmentation pathway, a frozen biomedical text
representation pathway, a shared transformation plan, and an evaluation-time
semantic audit (Fig.~\ref{fig:architecture}). The visual pathway fine-tunes the
CausalCLIPSeg RN50 encoder and retains its scale-wise supportive/inferior feature
partition, and two text-conditioned projection heads. Cached 512-dimensional
pooled BioMedCLIP features pass through a trainable state projection and condition
both heads; each evaluated head contains one active ATConv operator. During training, the shared plan
keeps the image and mask geometrically aligned and updates discrete spatial
phrases when their image-frame meaning changes. During evaluation, correct,
shuffled, fixed, and empty prompts test whether the checkpoint responds to
case-specific text. We use \emph{prompt-reliability-aware} in this operational
sense: the framework protects and audits prompt reliability. It does not estimate
per-case uncertainty or contain a learned reliability gate.

\begin{figure}[t]
\centering
\includegraphics[width=\linewidth]{medequiseg_architecture.pdf}
\caption{Architecture of \methodname{}. A public BUSI training case is used only
to illustrate the paired transformation. (A) During training, a keyed online plan
produces an aligned image--mask pair and the corresponding EquiPrompt lexical
view. (B) The fine-tuned CLIP RN50 features are partitioned into supportive and
inferior streams, where the latter is the complement of the supportive mask. A
trainable state projection maps the cached frozen BioMedCLIP feature
to $s_p$, which parameterizes the two causal projection heads; each head contains
one active ATConv operator. Their logits are averaged before segmentation loss.
The transformed target $y'$ provides the second input to the BCE--Dice loss.
Solid modules are optimized, dashed modules are frozen, and the gray feature is
cached offline. (C) At evaluation, true,
shuffled, fixed, and empty prompts audit semantic sensitivity of the same frozen
checkpoint; this audit is not a learned reliability gate.}
\label{fig:architecture}
\end{figure}

\subsubsection{Frozen follow-up shared plan}
The shared rows of the ordered ablation use the R11 follow-up policy defined
above. For case identifier $c$, epoch $e$, and training seed $s$, a keyed random
generator samples
\begin{equation}
  A_{c,e,s} = \mathcal{S}\!\left(H,V,R,F,Z,C,G,N\mid c,e,s\right),
\end{equation}
where $H$ and $V$ are horizontal and vertical flips, $R$ is a 90-degree
rotation, $F$ is a mild affine transform with translation, $Z$ is zoom, and
$C$, $G$, and $N$ denote color jitter, gamma, and Gaussian noise. The same
realized spatial plan is applied to image $x$ and mask $y$:
\begin{equation}
  x' = A^{x}_{c,e,s}(x), \qquad y' = A^{y}_{c,e,s}(y).
\end{equation}
Images use bilinear interpolation, masks use nearest-neighbor interpolation, and
the mask is never exposed to appearance operations. The training set is not
physically duplicated; a new deterministic plan is sampled online each epoch.

\begin{figure}[t]
\centering
\includegraphics[width=\linewidth]{equiprompt_mechanism.pdf}
\caption{EquiPrompt synchronization and its exact mechanism control, illustrated
with a public BUSI training case. (A) One
realized transform creates aligned image, mask, and prompt views; appearance-only
operations affect the image alone. (B) Direction phrases are represented in
image coordinates, transformed compositionally, and lexicalized only after all
discrete operations; multiword phrases are treated atomically. (C) \methodname{}
and \norewrite{} are separate runs matched in architecture, initialization,
optimizer, data order, and realized augmentation sequence. Only the lexical map
differs, isolating prompt validity from ordinary geometric augmentation without
assuming that the rewrite must improve accuracy.}
\label{fig:equiprompt}
\end{figure}

\subsubsection{Compositional direction rewriting}
Each recognized direction is mapped to a vector $d=(d_x,d_y)$ with components in
$\{-1,0,1\}$. Horizontal and vertical flips apply
\begin{equation}
  T_H(d_x,d_y)=(-d_x,d_y),\qquad
  T_V(d_x,d_y)=(d_x,-d_y),
\end{equation}
and a counter-clockwise 90-degree rotation applies
\begin{equation}
  T_R(d_x,d_y)=(d_y,-d_x).
\end{equation}
Sequential transforms compose these maps before converting the result back to
the original lexical family. Thus, \textit{upper right} remains an atomic target
rather than two unrelated replacements. Comma-separated targets are parsed as
separate atoms. The BioMedCLIP cache contains the finite lexical closure reachable
under all compositions, preventing an augmented prompt from silently falling
back to a missing embedding.

\subsubsection{Augmentation profile and scope}
\methodname{} uses one locked mild profile: $p_H=0.35$, $p_V=0.10$, and $p_R=0.12$;
when rotation is selected, $90^\circ$, $180^\circ$, and $270^\circ$ are sampled
uniformly. Mild affine is selected with probability 0.15, with angle in
$[-5^\circ,5^\circ]$ and translation in $\pm2.5\%$ of image size. Zoom is
selected with probability 0.25 in $[1.00,1.07]$. Color jitter, gamma, and noise
have probabilities 0.45, 0.25, and 0.25, respectively; jitter magnitude is 0.05,
gamma is in $[0.90,1.12]$, and noise standard deviation is 0.006 after
normalization.

Only flips and right-angle rotations trigger exact direction rewriting. Small
affine rotations, translation, and off-center zoom do not exchange coordinate
axes, but they can move a target across a coarse location bin. Consequently,
\methodname{} guarantees lexical consistency for the discrete axis-changing
operations, not for arbitrary continuous spatial descriptions. The small ranges
bound this approximation. We additionally ran \locrecompute{}, a benchmark-only
diagnostic that recomputes coarse position phrases from the connected-component
centroids of the transformed training mask whenever affine, translation, or zoom
is sampled. Although it uses the same transform distribution, model family, and
training seeds, its augmentation seed key differs from the main run and the
realized per-case transforms are therefore not exactly paired. We report it only
as an exploratory same-distribution diagnostic.

\subsubsection{Multimodal segmentation pathway}
The evaluated model initializes the CausalCLIPSeg visual pathway from the official
OpenAI CLIP RN50 checkpoint and fine-tunes it end to end. Prompts are represented
by cached, frozen 512-dimensional pooled features from
\texttt{microsoft/BiomedCLIP-PubMedBERT\_256-vit\_base\_patch16\_224}. A trainable linear projection
maps each pooled feature to the 1024-dimensional state consumed by both
CausalCLIPSeg projectors. Each projector converts this state into a sample-specific
$3\times3$ kernel and bias. One same-channel convolution in each active projector
is replaced by ATConv \cite{atconv}, giving two effective ATConv operators in the
evaluated graph. The adapter also instantiates expanded token features, but all
token-level fusion switches are disabled in the reported configuration; those
tokens therefore do not enter the effective forward path. Plugin metadata records
four substitutions because two additional replacements target dormant legacy FPN
members that are not called by the released \texttt{FPN\_AD.forward}; we exclude
them from the architectural claim. A post-run forward-hook audit loaded the
reported BUSI seed-123 checkpoint and passed one real test case through the
model: both projector operators were called once with non-zero outputs, while
the two registered neck candidates were never called (Supplementary Table S11). Training uses a BCE--Dice objective with Dice
weight 0.5. Text conditions visual decoding but does not directly generate the
target mask. \norewrite{} has the identical architecture, objective, and sampled
image--mask plans; only prompt rewriting is disabled.

\subsection{Experimental design}
\label{sec:experiments}
\subsubsection{Training and evaluation}
Inputs are resized to $224\times224$ for training. All \methodname{} runs use 100 epochs,
batch size 8, AdamW with initial learning rate $10^{-4}$ and weight decay
$10^{-4}$, and cosine decay to $10^{-6}$. Checkpoint selection uses validation
Dice. Evaluation uses a fixed threshold of 0.5. Predictions are restored from
the $224\times224$ inference grid to the native reference-mask grid before
scoring.
The principal \methodname{} and image-baseline environment used Python 3.11.15 and
PyTorch 2.5.1 with CUDA 12.1. The nnU-Net/MedCLIPSeg environment used
nnU-Net v2.8.0 and MONAI 1.3.0; the official U-Mamba path used Python 3.10.20
and nnU-Net v2.1.1. Runs used NVIDIA RTX 4090 GPUs with driver 580.119.02.
A package-level environment record lists the remaining core library versions.

Online complexity is measured from the frozen BUSI seed-123 checkpoints on one
NVIDIA GeForce RTX 4090 with batch size 1, FP32, and $224\times224$ inputs.
After model-specific warm-up, each forward pass is timed with wall-clock CUDA
synchronization; data loading and mask post-processing are excluded. The \methodname{}
and LViT measurements include cached-text lookup but not the one-time offline
BioMedCLIP or PubMedBERT encoding step. MedCLIPSeg instead runs its released
online UniMedCLIP/BiomedBERT path and is measured with both one probabilistic
sample and its official 30-sample test setting. Parameter counts include all
modules instantiated by each checkpoint loader and are not pruned to
output-connected subgraphs. Because the prompt-consistency policy is active only
during training, it adds no inference branch or parameter beyond the reported
\methodname{} backbone.

The expression \textit{same protocol} in this paper means the same locked
manifest, split roles, held-out test set, mask sanitation, prediction
restoration, and metric implementation. It does not imply that architectures
with different official recipes share identical losses, input sizes, or
augmentation. We therefore use the more precise phrase
\textit{same data and evaluation protocol} in comparative claims.

\subsubsection{Baselines and controls}
The completed comparison includes MedCLIPSeg, CausalCLIPSeg, LViT, RollingUNet,
U-Net, U-Net++, nnU-Net, U-Mamba, and U-KAN
\cite{medclipseg,causalclipseg,lvit,rollingunet,unet,unetpp,nnunet,umamba,ukan}.
The first three and \methodname{} are multimodal; the remaining six models are
image-only. Every method has three optimization seeds on every dataset. U-Net,
U-Net++, nnU-Net, U-Mamba, and U-KAN contribute 75 public
dataset--model--seed cells; RollingUNet contributes 15 more, for 90 image-only
public cells in total. The contextual CausalCLIPSeg row (internal run V3\_R0) and the
Base reference (V3\_ABL\_BASE) use the same CLIP-text/horizontal-flip recipe but
are independent three-seed training matrices; the latter is the reference for
the ordered component ablation, whereas the former is retained in the broad
model table.

\norewrite{} is the primary mechanism control. Four inference-time semantic controls are also
defined without retraining: correct case text, seed-fixed shuffled text,
a training-derived fixed prompt, and an empty prompt. Correct text must exceed
shuffled text in the Holm-adjusted paired rank screen before a dataset enters
the semantic-control visualization; a hierarchical interval excluding zero is
additionally required for the evidence to be described as robust to seed
uncertainty. The shuffled control applies a deterministic non-zero cyclic shift
(control seed 123) to the ordered test prompts. It preserves the complete prompt
multiset and assigns every row another row's prompt, although duplicated prompt
strings can remain textually unchanged; no class or attribute stratum is
preserved. To assess this confounding directly, a post hoc frozen-checkpoint
control permutes the complete prompt multiset separately within target-presence
and diagnostic-class strata. Strata are constructed only from released case
identifiers and prompt fields; the builder reads neither masks nor image pixels.
Within each stratum, a deterministic assignment maximizes changed prompt strings
while preserving one-to-one prompt closure. This changes 3,413 of 3,426 public
test prompts; 13 BUSI normal cases remain unchanged because every available
within-stratum prompt string is identical. Fixed and empty prompts are secondary
descriptive controls. The released target-derived captions on the
five public datasets define the formal but privileged referring benchmark.
Additional prompts generated directly from a case's mask for oracle diagnostic
experiments are a separate category and are excluded from formal performance
tables.

A frozen follow-up ordered component ablation uses the fixed V3 manifests and
evaluator for the base
CausalCLIPSeg pathway, BioMedCLIP substitution alone, ATConv alone, their
combination, the shared deterministic plan without rewriting, and the complete
\methodname{} configuration. Every configuration contains all five public
datasets and seeds 123, 456, and 789, yielding 90 unique
dataset--configuration--seed cells. The first four configurations use the same
hflip policy; the last two share the same richer online plan and differ only in
whether direction phrases are rewritten. Historical engineering runs using
other prompt-generation or evaluation definitions are excluded.

\subsubsection{Endpoints and statistical analysis}
The main accuracy-centered contrast reported in this retrospective ordered
ablation is the complete \methodname{} configuration minus the horizontal-flip Base
reference in equal-dataset Public-5 macro Dice. This contrast was selected after
the original frozen promotion gate and is therefore interpreted as post hoc
rather than as a prospectively declared primary endpoint. The mechanism
contrast frozen before completion of the paired matrix is \methodname{} minus \norewrite{} under an exactly matched
visual transformation plan. The ordered ablation additionally reports the
incremental BioMedCLIP and ATConv configurations. Both contrasts use
case-first paired inference after averaging the three seed values per case.
Dataset-specific Dice and IoU differences
and boundary measures are secondary endpoints. The primary semantic-control
endpoint is true versus shuffled Dice; fixed and empty controls are descriptive.
The presence/class-stratified derangement is a post hoc sensitivity analysis of
the true--shuffled contrast, not a replacement primary endpoint. It uses the
same case-first three-seed reduction, 10,000 paired bootstrap replicates, and
five-dataset Holm family as the other semantic comparisons.
Cross-model and cross-modality comparisons in Table~\ref{tab:bspc_v3_models}
are contextual rather than confirmatory because architectures and input
information differ.

Two additional training controls were completed after the 90-cell matrix. A
constant-text run keeps the complete recipe and shared plan but replaces every
cached prompt feature with the same empty-prompt embedding throughout training
and test. A Base-plus-plan run applies the richer shared plan to the Base-reference
recipe without BioMedCLIP substitution, ATConv, or rewriting. Their comparisons
with the earlier matrix are post hoc sensitivity analyses because the runner
code snapshot changed between batches, although manifests, evaluator, epochs,
optimizer settings, and command logs are retained. We therefore use these
controls to assess attribution, not as new confirmatory endpoints.

Seed-level results are reported as mean and sample standard deviation. For
hypothesis tests, per-case metrics are first averaged across seeds for each
method, then compared with paired Wilcoxon signed-rank tests; paired $t$-tests are
reported as sensitivity analyses. Family-wise multiplicity is controlled with
Holm correction \cite{holm}. Confidence intervals use 10,000 paired-case
bootstrap replicates after the three seed values for each case have been
averaged. For the Public-5 macro effect, resampling is stratified within each
dataset before the five dataset-level effects are averaged, so each benchmark
receives equal weight. Training-seed variability is retained in the reported
mean and sample standard deviation but is not treated as an additional set of
independent cases. Bootstrap intervals estimate the mean effect, whereas the
Wilcoxon test targets paired rank differences; the two summaries are interpreted
together rather than treated as interchangeable significance decisions.
Case-level tests use manifest rows as the paired units and are interpreted
subject to the patient-grouping limitation described below. Because the
COVID-19 test set contains repeated images from 474 recoverable subjects, a
post hoc sensitivity analysis retains the image-weighted point estimate but
resamples subjects and training seeds for the interval and applies paired tests
to equal-subject mean differences. This clustered analysis is reported
separately and does not replace the pre-specified image-level analysis.

\subsubsection{Software-assisted manuscript preparation}
OpenAI Codex was used for code-assisted reproducibility checks and English
language editing. The authors inspected the generated changes and retained full
responsibility for the experimental code, numerical results, citations, and
manuscript text.

\section{Results}
\label{sec:results}

\subsection{Segmentation accuracy and component contributions}
\def\BmcTableSelection{7}
% Tables for the BMC Medical Imaging manuscript.
% Ordered by the evidence hierarchy: dataset audit, ordered configurations,
% semantic controls, contextual model comparison, and exploratory diagnostics.

\ifnum\BmcTableSelection=1
\begin{table}[htbp]
\centering
\caption{Image-level dataset partitions and realized prompt-change frequencies.}
\label{tab:bspc_dataset_audit}
\scriptsize
\setlength{\tabcolsep}{2.0pt}
\begin{tabular}{@{}lrrrlrrr@{}}
\toprule
Dataset & Train & Val. & Test & Prompt & Unique & Spatial (\%) & Change (\%) \\
\midrule
BUSI     & 624  & 78  & 78   & Target-derived & 198 & 82.85 & $20.27\pm0.10$ \\
ClinicDB & 490  & 61  & 61   & Target-derived & 202 & 100.00 & $28.48\pm0.21$ \\
BUS-BRA  & 1311 & 282 & 282  & Target-derived & 75  & 100.00 & $12.44\pm0.12$ \\
BRISC    & 3881 & 974 & 892  & Target-derived & 359 & 78.23 & $5.75\pm0.04$ \\
COVID-19 & 5716 & 1429& 2113 & Target-derived & 269 & 100.00 & $45.75\pm0.06$ \\
\bottomrule
\end{tabular}
\vspace{2pt}
\begin{minipage}{0.98\linewidth}
\footnotesize\raggedright
\textit{Notes:} Unique and spatial coverage use training prompts. Realized
change enumerates all case--epoch events under the fixed 100-epoch shared
augmentation policy and reports mean$\pm$sample standard deviation across seeds
123, 456, and 789. Public prompts are target-derived synthetic captions, not
clinical text.\par
\end{minipage}
\end{table}
\fi

\ifnum\BmcTableSelection=3
\begin{table}[htbp]
\centering
\caption{Inference-time prompt controls for trained \methodname{} models on
five public datasets.}
\label{tab:bspc_semantic_controls}
\scriptsize
\setlength{\tabcolsep}{4.5pt}
\begin{tabular}{@{}lrrrrr@{}}
\toprule
Dataset & True Dice & Shuffled Dice & Fixed Dice & Empty Dice & $\Delta$ pp [95\% CI] \\
\midrule
BUSI     & $85.78\pm0.45$ & $46.73\pm1.41$ & $16.69\pm0.02$ & $13.22\pm5.35$ & $+39.05$ [$29.85,48.49$] \\
ClinicDB & $92.07\pm0.71$ & $90.20\pm0.94$ & $88.31\pm1.31$ & $78.71\pm7.93$ & $+1.88$ [$0.07,4.22$] \\
BUS-BRA  & $90.24\pm0.14$ & $89.47\pm0.23$ & $88.87\pm0.28$ & $68.22\pm8.44$ & $+0.77$ [$0.45,1.15$] \\
BRISC    & $87.18\pm0.14$ & $70.89\pm2.12$ & $75.31\pm1.97$ & $58.89\pm5.30$ & $+16.29$ [$14.26,18.43$] \\
COVID-19 & $84.08\pm0.13$ & $68.80\pm0.82$ & $68.54\pm0.16$ & $55.81\pm3.82$ & $+15.29$ [$14.49,16.10$] \\
\bottomrule
\end{tabular}
\vspace{2pt}
\begin{minipage}{0.98\linewidth}
\footnotesize\raggedright
\textit{Notes:} No model is retrained; only the test prompt changes. Dice is
mean$\pm$sample standard deviation over three seeds. The last column is true
minus shuffled Dice with a paired-case bootstrap 95\% interval after
three-seed averaging. Holm-adjusted Wilcoxon tests and the bootstrap intervals
are positive on all five datasets.\par
\end{minipage}
\end{table}
\fi

\ifnum\BmcTableSelection=4
% Auto-generated by generate_medequiseg_factorial_manuscript_assets.py.
% Do not edit numerical cells manually.
\begin{table}[htbp]
\centering
\caption{Five-dataset macro-average accuracy across five retained ordered configurations.}
\label{tab:medequiseg_factorial}
\scriptsize
\setlength{\tabcolsep}{2.4pt}
\begin{tabular}{@{}lccccrr@{}}
\toprule
Configuration & Text enc. & ATConv & Aug. & Rewrite & Dice & IoU \\
\midrule
Base model & -- & -- & -- & -- & $86.41\pm0.35$ & $79.22\pm0.34$ \\
$+$ BioMedCLIP & \checkmark & -- & -- & -- & $86.67\pm0.08$ & $79.52\pm0.06$ \\
$+$ ATConv & -- & \checkmark & -- & -- & $86.63\pm0.32$ & $79.42\pm0.31$ \\
$+$ BioMedCLIP $+$ ATConv & \checkmark & \checkmark & -- & -- & \underline{$87.03\pm0.15$} & \underline{$79.85\pm0.23$} \\
\textbf{\methodname{} (Ours)} & \checkmark & \checkmark & \checkmark & \checkmark & \textbf{$87.87\pm0.17$} & \textbf{$80.97\pm0.13$} \\
\bottomrule
\end{tabular}
\vspace{2pt}
\begin{minipage}{0.98\linewidth}
\footnotesize\raggedright
\textit{Notes:} Five retained configurations, five datasets, and three seeds produce 75
training runs. Values are five-dataset macro-average percentages
(mean$\pm$sample standard deviation). Text enc. denotes pooled BioMedCLIP
conditioning; ATConv denotes the two projection-head operators active in the evaluated forward pass;
Aug. denotes the common shared augmentation policy. Shared augmentation and
prompt rewriting appear together only in the final configuration, so these are conditional configuration contrasts,
not independent component effects or complete interactions. Bold and underline
mark the best and second-best configurations. The complete per-dataset matrix
is in Supplementary Table S9.\par
\end{minipage}
\end{table}

\fi

\ifnum\BmcTableSelection=5
\begin{table}[htbp]
\centering
\caption{Per-dataset and equal-dataset macro Dice on five public benchmarks.}
\label{tab:bspc_v3_models}
\scriptsize
\setlength{\tabcolsep}{2.0pt}
\begin{tabular}{@{}lrrrrrr@{}}
\toprule
Method & BUSI & ClinicDB & BUS-BRA & BRISC & COVID-19 & Macro avg. \\
\midrule
\multicolumn{7}{l}{\textit{Image+text methods}} \\
MedCLIPSeg
 & $\mathbf{86.13}$ & $\underline{91.96}$ & $88.47$
 & $81.36$ & $\underline{82.62}$ & $86.11$ \\
CausalCLIPSeg
 & $84.62$ & $91.19$ & $\underline{88.87}$
 & $\underline{86.11}$ & $81.85$ & $\underline{86.53}$ \\
LViT
 & $81.81$ & $91.03$ & $87.66$
 & $85.24$ & $78.50$ & $84.85$ \\
\textbf{\methodname{} (Ours)}
 & $\underline{85.78}$ & $\mathbf{92.07}$ & $\mathbf{90.24}$
 & $\mathbf{87.18}$ & $\mathbf{84.08}$ & $\mathbf{87.87}$ \\
\midrule
\multicolumn{7}{l}{\textit{Image-only contextual references}} \\
RollingUNet
 & $72.36$ & $87.31$ & $88.00$
 & $85.03$ & $79.32$ & $82.40$ \\
U-Net
 & $61.31$ & $85.61$ & $87.35$
 & $85.00$ & $78.81$ & $79.62$ \\
U-Net++
 & $76.21$ & $89.77$ & $88.59$
 & $\underline{85.94}$ & $\mathbf{80.00}$ & $84.10$ \\
nnU-Net
 & $74.99$ & $\mathbf{93.62}$ & $\underline{89.80}$
 & $\mathbf{87.30}$ & $79.07$ & $\mathbf{84.96}$ \\
U-Mamba
 & $\underline{76.63}$ & $\underline{92.21}$ & $\mathbf{90.34}$
 & $85.37$ & $79.08$ & $\underline{84.73}$ \\
U-KAN
 & $\mathbf{77.10}$ & $91.40$ & $88.56$
 & $85.45$ & $\underline{79.38}$ & $84.38$ \\
\bottomrule
\end{tabular}
\vspace{2pt}
\begin{minipage}{0.98\linewidth}
\footnotesize\raggedright
\textit{Notes:} Dataset values are three-seed mean Dice percentages; the macro
average is the arithmetic mean of the five dataset means. Bold and underline rank
methods only within the image+text
or image-only panel. Cross-panel differences are descriptive because the
public image+text methods receive target-derived synthetic captions unavailable
to image-only methods. Complete mean$\pm$sample standard deviation Dice and IoU
results are in Supplementary Table S14. The five \methodname{} entries use the
models from the ordered configuration analysis.\par
\end{minipage}
\end{table}
\fi

\ifnum\BmcTableSelection=7
\begin{table}[t]
\centering
\caption{Equal-dataset macro segmentation accuracy on the five public
benchmarks. Values are the arithmetic mean of the three-seed dataset means in
Table~\ref{tab:bspc_v3_models}. Bold and underline rank
models within modality. Cross-modality differences are descriptive because the
image+text models receive target-derived synthetic captions unavailable to the
image-only models.}
\label{tab:public5_macro_accuracy}
\small
\setlength{\tabcolsep}{6pt}
\begin{tabular}{lcc}
\toprule
Method & Dice (\%) & IoU (\%) \\
\midrule
\multicolumn{3}{l}{\textit{Image+text methods}} \\
\textbf{\methodname{} (Ours)} & $\mathbf{87.87}$ & $\mathbf{80.97}$ \\
CausalCLIPSeg Baseline & $\underline{86.53}$ & $\underline{79.32}$ \\
MedCLIPSeg & $86.11$ & $78.27$ \\
LViT & $84.85$ & $77.37$ \\
\midrule
\multicolumn{3}{l}{\textit{Image-only contextual references}} \\
nnU-Net & $\mathbf{84.96}$ & $\mathbf{78.23}$ \\
U-Mamba & $\underline{84.73}$ & $\underline{77.94}$ \\
U-KAN & $84.38$ & $77.02$ \\
U-Net++ & $84.10$ & $76.65$ \\
RollingUNet & $82.40$ & $74.62$ \\
U-Net & $79.62$ & $71.20$ \\
\bottomrule
\end{tabular}
\end{table}
\fi

\ifnum\BmcTableSelection=6
\begin{table}[htbp]
\centering
\caption{Exploratory three-seed continuous-location recomputation analysis.
\locrecompute{}
recomputes coarse position phrases from the transformed training mask when a
mild affine, translation, or zoom is sampled. Values are mean$\pm$sample
standard deviation. Delta and the reported bootstrap interval are
\locrecompute{} minus its reference model in percentage points. Both runs use
the same transformation distribution but independently initialized augmentation
sequences, so the realized per-case transformations are not paired. These
rows are not pooled with the main ordered configuration analysis. $p_W$ is the exploratory, unadjusted two-sided
paired Wilcoxon $p$-value.}
\label{tab:continuous_location_diagnostic}
\small
\setlength{\tabcolsep}{4.2pt}
\begin{tabular}{@{}llrrrr@{}}
\toprule
Dataset & Metric & Reference model & Location Recompute & $\Delta$ pp [95\% CI] & $p_W$ \\
\midrule
BUSI & Dice & $85.74\pm0.80$ & $84.62\pm0.20$ & $-1.12$ [$-2.62,0.07$] & .004 \\
     & IoU  & $79.46\pm0.83$ & $78.06\pm0.64$ & $-1.40$ [$-2.70,-0.30$] & .002 \\
ClinicDB & Dice & $93.18\pm0.54$ & $93.06\pm0.72$ & $-0.13$ [$-1.25,1.27$] & .138 \\
         & IoU  & $87.77\pm0.82$ & $87.54\pm1.07$ & $-0.23$ [$-1.94,1.78$] & .117 \\
\bottomrule
\end{tabular}
\end{table}
\fi


Across the five public datasets, \methodname{} attains equal-dataset macro Dice
of $87.87\%$ and IoU of $80.97\%$. This is $1.34$ Dice pp and $1.65$ IoU pp
above the next reproduced image+text model, CausalCLIPSeg, and $2.91$ Dice pp
and $2.74$ IoU pp above the highest-macro image-only architecture, nnU-Net.
These macro comparisons summarize the fixed public benchmark panel; the
cross-modality difference remains descriptive because image+text methods receive
target-derived captions. Per-dataset results are reported below to expose
heterogeneity rather than treating the macro value as a universal ranking.

\def\BmcTableSelection{4}
% Tables for the BMC Medical Imaging manuscript.
% Ordered by the evidence hierarchy: dataset audit, ordered configurations,
% semantic controls, contextual model comparison, and exploratory diagnostics.

\ifnum\BmcTableSelection=1
\begin{table}[htbp]
\centering
\caption{Image-level dataset partitions and realized prompt-change frequencies.}
\label{tab:bspc_dataset_audit}
\scriptsize
\setlength{\tabcolsep}{2.0pt}
\begin{tabular}{@{}lrrrlrrr@{}}
\toprule
Dataset & Train & Val. & Test & Prompt & Unique & Spatial (\%) & Change (\%) \\
\midrule
BUSI     & 624  & 78  & 78   & Target-derived & 198 & 82.85 & $20.27\pm0.10$ \\
ClinicDB & 490  & 61  & 61   & Target-derived & 202 & 100.00 & $28.48\pm0.21$ \\
BUS-BRA  & 1311 & 282 & 282  & Target-derived & 75  & 100.00 & $12.44\pm0.12$ \\
BRISC    & 3881 & 974 & 892  & Target-derived & 359 & 78.23 & $5.75\pm0.04$ \\
COVID-19 & 5716 & 1429& 2113 & Target-derived & 269 & 100.00 & $45.75\pm0.06$ \\
\bottomrule
\end{tabular}
\vspace{2pt}
\begin{minipage}{0.98\linewidth}
\footnotesize\raggedright
\textit{Notes:} Unique and spatial coverage use training prompts. Realized
change enumerates all case--epoch events under the fixed 100-epoch shared
augmentation policy and reports mean$\pm$sample standard deviation across seeds
123, 456, and 789. Public prompts are target-derived synthetic captions, not
clinical text.\par
\end{minipage}
\end{table}
\fi

\ifnum\BmcTableSelection=3
\begin{table}[htbp]
\centering
\caption{Inference-time prompt controls for trained \methodname{} models on
five public datasets.}
\label{tab:bspc_semantic_controls}
\scriptsize
\setlength{\tabcolsep}{4.5pt}
\begin{tabular}{@{}lrrrrr@{}}
\toprule
Dataset & True Dice & Shuffled Dice & Fixed Dice & Empty Dice & $\Delta$ pp [95\% CI] \\
\midrule
BUSI     & $85.78\pm0.45$ & $46.73\pm1.41$ & $16.69\pm0.02$ & $13.22\pm5.35$ & $+39.05$ [$29.85,48.49$] \\
ClinicDB & $92.07\pm0.71$ & $90.20\pm0.94$ & $88.31\pm1.31$ & $78.71\pm7.93$ & $+1.88$ [$0.07,4.22$] \\
BUS-BRA  & $90.24\pm0.14$ & $89.47\pm0.23$ & $88.87\pm0.28$ & $68.22\pm8.44$ & $+0.77$ [$0.45,1.15$] \\
BRISC    & $87.18\pm0.14$ & $70.89\pm2.12$ & $75.31\pm1.97$ & $58.89\pm5.30$ & $+16.29$ [$14.26,18.43$] \\
COVID-19 & $84.08\pm0.13$ & $68.80\pm0.82$ & $68.54\pm0.16$ & $55.81\pm3.82$ & $+15.29$ [$14.49,16.10$] \\
\bottomrule
\end{tabular}
\vspace{2pt}
\begin{minipage}{0.98\linewidth}
\footnotesize\raggedright
\textit{Notes:} No model is retrained; only the test prompt changes. Dice is
mean$\pm$sample standard deviation over three seeds. The last column is true
minus shuffled Dice with a paired-case bootstrap 95\% interval after
three-seed averaging. Holm-adjusted Wilcoxon tests and the bootstrap intervals
are positive on all five datasets.\par
\end{minipage}
\end{table}
\fi

\ifnum\BmcTableSelection=4
% Auto-generated by generate_medequiseg_factorial_manuscript_assets.py.
% Do not edit numerical cells manually.
\begin{table}[htbp]
\centering
\caption{Five-dataset macro-average accuracy across five retained ordered configurations.}
\label{tab:medequiseg_factorial}
\scriptsize
\setlength{\tabcolsep}{2.4pt}
\begin{tabular}{@{}lccccrr@{}}
\toprule
Configuration & Text enc. & ATConv & Aug. & Rewrite & Dice & IoU \\
\midrule
Base model & -- & -- & -- & -- & $86.41\pm0.35$ & $79.22\pm0.34$ \\
$+$ BioMedCLIP & \checkmark & -- & -- & -- & $86.67\pm0.08$ & $79.52\pm0.06$ \\
$+$ ATConv & -- & \checkmark & -- & -- & $86.63\pm0.32$ & $79.42\pm0.31$ \\
$+$ BioMedCLIP $+$ ATConv & \checkmark & \checkmark & -- & -- & \underline{$87.03\pm0.15$} & \underline{$79.85\pm0.23$} \\
\textbf{\methodname{} (Ours)} & \checkmark & \checkmark & \checkmark & \checkmark & \textbf{$87.87\pm0.17$} & \textbf{$80.97\pm0.13$} \\
\bottomrule
\end{tabular}
\vspace{2pt}
\begin{minipage}{0.98\linewidth}
\footnotesize\raggedright
\textit{Notes:} Five retained configurations, five datasets, and three seeds produce 75
training runs. Values are five-dataset macro-average percentages
(mean$\pm$sample standard deviation). Text enc. denotes pooled BioMedCLIP
conditioning; ATConv denotes the two projection-head operators active in the evaluated forward pass;
Aug. denotes the common shared augmentation policy. Shared augmentation and
prompt rewriting appear together only in the final configuration, so these are conditional configuration contrasts,
not independent component effects or complete interactions. Bold and underline
mark the best and second-best configurations. The complete per-dataset matrix
is in Supplementary Table S9.\par
\end{minipage}
\end{table}

\fi

\ifnum\BmcTableSelection=5
\begin{table}[htbp]
\centering
\caption{Per-dataset and equal-dataset macro Dice on five public benchmarks.}
\label{tab:bspc_v3_models}
\scriptsize
\setlength{\tabcolsep}{2.0pt}
\begin{tabular}{@{}lrrrrrr@{}}
\toprule
Method & BUSI & ClinicDB & BUS-BRA & BRISC & COVID-19 & Macro avg. \\
\midrule
\multicolumn{7}{l}{\textit{Image+text methods}} \\
MedCLIPSeg
 & $\mathbf{86.13}$ & $\underline{91.96}$ & $88.47$
 & $81.36$ & $\underline{82.62}$ & $86.11$ \\
CausalCLIPSeg
 & $84.62$ & $91.19$ & $\underline{88.87}$
 & $\underline{86.11}$ & $81.85$ & $\underline{86.53}$ \\
LViT
 & $81.81$ & $91.03$ & $87.66$
 & $85.24$ & $78.50$ & $84.85$ \\
\textbf{\methodname{} (Ours)}
 & $\underline{85.78}$ & $\mathbf{92.07}$ & $\mathbf{90.24}$
 & $\mathbf{87.18}$ & $\mathbf{84.08}$ & $\mathbf{87.87}$ \\
\midrule
\multicolumn{7}{l}{\textit{Image-only contextual references}} \\
RollingUNet
 & $72.36$ & $87.31$ & $88.00$
 & $85.03$ & $79.32$ & $82.40$ \\
U-Net
 & $61.31$ & $85.61$ & $87.35$
 & $85.00$ & $78.81$ & $79.62$ \\
U-Net++
 & $76.21$ & $89.77$ & $88.59$
 & $\underline{85.94}$ & $\mathbf{80.00}$ & $84.10$ \\
nnU-Net
 & $74.99$ & $\mathbf{93.62}$ & $\underline{89.80}$
 & $\mathbf{87.30}$ & $79.07$ & $\mathbf{84.96}$ \\
U-Mamba
 & $\underline{76.63}$ & $\underline{92.21}$ & $\mathbf{90.34}$
 & $85.37$ & $79.08$ & $\underline{84.73}$ \\
U-KAN
 & $\mathbf{77.10}$ & $91.40$ & $88.56$
 & $85.45$ & $\underline{79.38}$ & $84.38$ \\
\bottomrule
\end{tabular}
\vspace{2pt}
\begin{minipage}{0.98\linewidth}
\footnotesize\raggedright
\textit{Notes:} Dataset values are three-seed mean Dice percentages; the macro
average is the arithmetic mean of the five dataset means. Bold and underline rank
methods only within the image+text
or image-only panel. Cross-panel differences are descriptive because the
public image+text methods receive target-derived synthetic captions unavailable
to image-only methods. Complete mean$\pm$sample standard deviation Dice and IoU
results are in Supplementary Table S14. The five \methodname{} entries use the
models from the ordered configuration analysis.\par
\end{minipage}
\end{table}
\fi

\ifnum\BmcTableSelection=7
\begin{table}[t]
\centering
\caption{Equal-dataset macro segmentation accuracy on the five public
benchmarks. Values are the arithmetic mean of the three-seed dataset means in
Table~\ref{tab:bspc_v3_models}. Bold and underline rank
models within modality. Cross-modality differences are descriptive because the
image+text models receive target-derived synthetic captions unavailable to the
image-only models.}
\label{tab:public5_macro_accuracy}
\small
\setlength{\tabcolsep}{6pt}
\begin{tabular}{lcc}
\toprule
Method & Dice (\%) & IoU (\%) \\
\midrule
\multicolumn{3}{l}{\textit{Image+text methods}} \\
\textbf{\methodname{} (Ours)} & $\mathbf{87.87}$ & $\mathbf{80.97}$ \\
CausalCLIPSeg Baseline & $\underline{86.53}$ & $\underline{79.32}$ \\
MedCLIPSeg & $86.11$ & $78.27$ \\
LViT & $84.85$ & $77.37$ \\
\midrule
\multicolumn{3}{l}{\textit{Image-only contextual references}} \\
nnU-Net & $\mathbf{84.96}$ & $\mathbf{78.23}$ \\
U-Mamba & $\underline{84.73}$ & $\underline{77.94}$ \\
U-KAN & $84.38$ & $77.02$ \\
U-Net++ & $84.10$ & $76.65$ \\
RollingUNet & $82.40$ & $74.62$ \\
U-Net & $79.62$ & $71.20$ \\
\bottomrule
\end{tabular}
\end{table}
\fi

\ifnum\BmcTableSelection=6
\begin{table}[htbp]
\centering
\caption{Exploratory three-seed continuous-location recomputation analysis.
\locrecompute{}
recomputes coarse position phrases from the transformed training mask when a
mild affine, translation, or zoom is sampled. Values are mean$\pm$sample
standard deviation. Delta and the reported bootstrap interval are
\locrecompute{} minus its reference model in percentage points. Both runs use
the same transformation distribution but independently initialized augmentation
sequences, so the realized per-case transformations are not paired. These
rows are not pooled with the main ordered configuration analysis. $p_W$ is the exploratory, unadjusted two-sided
paired Wilcoxon $p$-value.}
\label{tab:continuous_location_diagnostic}
\small
\setlength{\tabcolsep}{4.2pt}
\begin{tabular}{@{}llrrrr@{}}
\toprule
Dataset & Metric & Reference model & Location Recompute & $\Delta$ pp [95\% CI] & $p_W$ \\
\midrule
BUSI & Dice & $85.74\pm0.80$ & $84.62\pm0.20$ & $-1.12$ [$-2.62,0.07$] & .004 \\
     & IoU  & $79.46\pm0.83$ & $78.06\pm0.64$ & $-1.40$ [$-2.70,-0.30$] & .002 \\
ClinicDB & Dice & $93.18\pm0.54$ & $93.06\pm0.72$ & $-0.13$ [$-1.25,1.27$] & .138 \\
         & IoU  & $87.77\pm0.82$ & $87.54\pm1.07$ & $-0.23$ [$-1.94,1.78$] & .117 \\
\bottomrule
\end{tabular}
\end{table}
\fi


The frozen follow-up ordered component ablation contains all 90 planned cells. Public-5 macro Dice rises
from $86.41\pm0.35\%$ for the base pathway to $86.67\pm0.08\%$ with BioMedCLIP,
$86.63\pm0.32\%$ with ATConv, and $87.03\pm0.15\%$ with both components. Adding
the shared deterministic plan without rewriting gives $87.96\pm0.05\%$; the
complete \methodname{} configuration gives $87.87\pm0.17\%$. The corresponding
full-framework minus base-model differences are $+1.46$ Dice pp (95\%
dataset-stratified bootstrap CI $+0.95$ to $+1.99$) and $+1.74$ IoU pp (95\% CI
$+1.13$ to $+2.34$). Dataset-level Dice point estimates are positive throughout;
only the ClinicDB interval includes zero. The component increments are
descriptive because this ordered ablation does
not estimate every higher-order interaction, but it shows that neither
BioMedCLIP nor ATConv alone accounts for the final improvement.

\begin{table}[htbp]
\centering
\caption{Selected contrast from the 75-run ordered configuration
study. Deltas are percentage points with 95\% dataset-stratified paired-case
bootstrap confidence intervals after averaging each case across the three
training seeds. These contrasts describe differences between complete
configurations and do not identify independent causal effects of individual
components.}
\label{tab:attribution_controls}
\scriptsize
\setlength{\tabcolsep}{3.5pt}
\begin{tabular}{@{}lrr@{}}
\toprule
Comparison & Dice $\Delta$ [95\% CI] & IoU $\Delta$ [95\% CI] \\
\midrule
\methodname{} $-$ Base reference & $+1.46$ [$+0.95,+1.99$] & $+1.74$ [$+1.13,+2.34$] \\
\bottomrule
\end{tabular}
\end{table}


The follow-up controls further localize the aggregate effect. Base plus the
shared plan reaches $87.83\pm0.23\%$ Dice and differs from the Base reference by $+1.42$
pp (95\% CI $+0.81$ to $+2.13$). The complete model differs from Base plus the
shared plan by only $+0.04$ pp (95\% CI $-0.40$ to $+0.41$), providing no
aggregate evidence for an additional contribution from the combined
BioMedCLIP--ATConv--rewrite stack beyond the plan. Conversely, the complete
model differs from its constant-text training control by $+1.90$ pp (95\% CI
$+1.16$ to $+2.75$), with intervals excluding zero on BUSI and COVID-19 but not
on ClinicDB, BUS-BRA, or BRISC. Thus case-specific privileged text matters in
aggregate, but its evidence is concentrated rather than uniform.

\begin{figure}[t]
\centering
\includegraphics[width=\linewidth]{strict_factorial_evidence.pdf}
\caption{Accuracy-centered follow-up ordered component ablation on the fixed V3 data and evaluation protocol.
(A) Public-5 macro Dice for all six configurations; points are three-seed means
and bars are sample standard deviations. (B) Paired three-seed Dice differences
between \methodname{} and the Base pathway for each public dataset and their
equal-dataset macro; bars are sample standard deviations. (C) Dice differences
between \methodname{} and the shared-plan no-rewrite control. Dataset bars and
the macro diamond are 95\% paired-case bootstrap intervals after three-seed
case averaging. The Public-5 macro difference from Base is positive,
whereas the rewrite-only contrast is centered near zero.}
\label{fig:strict_factorial}
\end{figure}
\FloatBarrier

\subsubsection{Matched rewrite and semantic-control evidence}
\def\BmcTableSelection{2}
% Tables for the BMC Medical Imaging manuscript.
% Ordered by the evidence hierarchy: dataset audit, ordered configurations,
% semantic controls, contextual model comparison, and exploratory diagnostics.

\ifnum\BmcTableSelection=1
\begin{table}[htbp]
\centering
\caption{Image-level dataset partitions and realized prompt-change frequencies.}
\label{tab:bspc_dataset_audit}
\scriptsize
\setlength{\tabcolsep}{2.0pt}
\begin{tabular}{@{}lrrrlrrr@{}}
\toprule
Dataset & Train & Val. & Test & Prompt & Unique & Spatial (\%) & Change (\%) \\
\midrule
BUSI     & 624  & 78  & 78   & Target-derived & 198 & 82.85 & $20.27\pm0.10$ \\
ClinicDB & 490  & 61  & 61   & Target-derived & 202 & 100.00 & $28.48\pm0.21$ \\
BUS-BRA  & 1311 & 282 & 282  & Target-derived & 75  & 100.00 & $12.44\pm0.12$ \\
BRISC    & 3881 & 974 & 892  & Target-derived & 359 & 78.23 & $5.75\pm0.04$ \\
COVID-19 & 5716 & 1429& 2113 & Target-derived & 269 & 100.00 & $45.75\pm0.06$ \\
\bottomrule
\end{tabular}
\vspace{2pt}
\begin{minipage}{0.98\linewidth}
\footnotesize\raggedright
\textit{Notes:} Unique and spatial coverage use training prompts. Realized
change enumerates all case--epoch events under the fixed 100-epoch shared
augmentation policy and reports mean$\pm$sample standard deviation across seeds
123, 456, and 789. Public prompts are target-derived synthetic captions, not
clinical text.\par
\end{minipage}
\end{table}
\fi

\ifnum\BmcTableSelection=3
\begin{table}[htbp]
\centering
\caption{Inference-time prompt controls for trained \methodname{} models on
five public datasets.}
\label{tab:bspc_semantic_controls}
\scriptsize
\setlength{\tabcolsep}{4.5pt}
\begin{tabular}{@{}lrrrrr@{}}
\toprule
Dataset & True Dice & Shuffled Dice & Fixed Dice & Empty Dice & $\Delta$ pp [95\% CI] \\
\midrule
BUSI     & $85.78\pm0.45$ & $46.73\pm1.41$ & $16.69\pm0.02$ & $13.22\pm5.35$ & $+39.05$ [$29.85,48.49$] \\
ClinicDB & $92.07\pm0.71$ & $90.20\pm0.94$ & $88.31\pm1.31$ & $78.71\pm7.93$ & $+1.88$ [$0.07,4.22$] \\
BUS-BRA  & $90.24\pm0.14$ & $89.47\pm0.23$ & $88.87\pm0.28$ & $68.22\pm8.44$ & $+0.77$ [$0.45,1.15$] \\
BRISC    & $87.18\pm0.14$ & $70.89\pm2.12$ & $75.31\pm1.97$ & $58.89\pm5.30$ & $+16.29$ [$14.26,18.43$] \\
COVID-19 & $84.08\pm0.13$ & $68.80\pm0.82$ & $68.54\pm0.16$ & $55.81\pm3.82$ & $+15.29$ [$14.49,16.10$] \\
\bottomrule
\end{tabular}
\vspace{2pt}
\begin{minipage}{0.98\linewidth}
\footnotesize\raggedright
\textit{Notes:} No model is retrained; only the test prompt changes. Dice is
mean$\pm$sample standard deviation over three seeds. The last column is true
minus shuffled Dice with a paired-case bootstrap 95\% interval after
three-seed averaging. Holm-adjusted Wilcoxon tests and the bootstrap intervals
are positive on all five datasets.\par
\end{minipage}
\end{table}
\fi

\ifnum\BmcTableSelection=4
% Auto-generated by generate_medequiseg_factorial_manuscript_assets.py.
% Do not edit numerical cells manually.
\begin{table}[htbp]
\centering
\caption{Five-dataset macro-average accuracy across five retained ordered configurations.}
\label{tab:medequiseg_factorial}
\scriptsize
\setlength{\tabcolsep}{2.4pt}
\begin{tabular}{@{}lccccrr@{}}
\toprule
Configuration & Text enc. & ATConv & Aug. & Rewrite & Dice & IoU \\
\midrule
Base model & -- & -- & -- & -- & $86.41\pm0.35$ & $79.22\pm0.34$ \\
$+$ BioMedCLIP & \checkmark & -- & -- & -- & $86.67\pm0.08$ & $79.52\pm0.06$ \\
$+$ ATConv & -- & \checkmark & -- & -- & $86.63\pm0.32$ & $79.42\pm0.31$ \\
$+$ BioMedCLIP $+$ ATConv & \checkmark & \checkmark & -- & -- & \underline{$87.03\pm0.15$} & \underline{$79.85\pm0.23$} \\
\textbf{\methodname{} (Ours)} & \checkmark & \checkmark & \checkmark & \checkmark & \textbf{$87.87\pm0.17$} & \textbf{$80.97\pm0.13$} \\
\bottomrule
\end{tabular}
\vspace{2pt}
\begin{minipage}{0.98\linewidth}
\footnotesize\raggedright
\textit{Notes:} Five retained configurations, five datasets, and three seeds produce 75
training runs. Values are five-dataset macro-average percentages
(mean$\pm$sample standard deviation). Text enc. denotes pooled BioMedCLIP
conditioning; ATConv denotes the two projection-head operators active in the evaluated forward pass;
Aug. denotes the common shared augmentation policy. Shared augmentation and
prompt rewriting appear together only in the final configuration, so these are conditional configuration contrasts,
not independent component effects or complete interactions. Bold and underline
mark the best and second-best configurations. The complete per-dataset matrix
is in Supplementary Table S9.\par
\end{minipage}
\end{table}

\fi

\ifnum\BmcTableSelection=5
\begin{table}[htbp]
\centering
\caption{Per-dataset and equal-dataset macro Dice on five public benchmarks.}
\label{tab:bspc_v3_models}
\scriptsize
\setlength{\tabcolsep}{2.0pt}
\begin{tabular}{@{}lrrrrrr@{}}
\toprule
Method & BUSI & ClinicDB & BUS-BRA & BRISC & COVID-19 & Macro avg. \\
\midrule
\multicolumn{7}{l}{\textit{Image+text methods}} \\
MedCLIPSeg
 & $\mathbf{86.13}$ & $\underline{91.96}$ & $88.47$
 & $81.36$ & $\underline{82.62}$ & $86.11$ \\
CausalCLIPSeg
 & $84.62$ & $91.19$ & $\underline{88.87}$
 & $\underline{86.11}$ & $81.85$ & $\underline{86.53}$ \\
LViT
 & $81.81$ & $91.03$ & $87.66$
 & $85.24$ & $78.50$ & $84.85$ \\
\textbf{\methodname{} (Ours)}
 & $\underline{85.78}$ & $\mathbf{92.07}$ & $\mathbf{90.24}$
 & $\mathbf{87.18}$ & $\mathbf{84.08}$ & $\mathbf{87.87}$ \\
\midrule
\multicolumn{7}{l}{\textit{Image-only contextual references}} \\
RollingUNet
 & $72.36$ & $87.31$ & $88.00$
 & $85.03$ & $79.32$ & $82.40$ \\
U-Net
 & $61.31$ & $85.61$ & $87.35$
 & $85.00$ & $78.81$ & $79.62$ \\
U-Net++
 & $76.21$ & $89.77$ & $88.59$
 & $\underline{85.94}$ & $\mathbf{80.00}$ & $84.10$ \\
nnU-Net
 & $74.99$ & $\mathbf{93.62}$ & $\underline{89.80}$
 & $\mathbf{87.30}$ & $79.07$ & $\mathbf{84.96}$ \\
U-Mamba
 & $\underline{76.63}$ & $\underline{92.21}$ & $\mathbf{90.34}$
 & $85.37$ & $79.08$ & $\underline{84.73}$ \\
U-KAN
 & $\mathbf{77.10}$ & $91.40$ & $88.56$
 & $85.45$ & $\underline{79.38}$ & $84.38$ \\
\bottomrule
\end{tabular}
\vspace{2pt}
\begin{minipage}{0.98\linewidth}
\footnotesize\raggedright
\textit{Notes:} Dataset values are three-seed mean Dice percentages; the macro
average is the arithmetic mean of the five dataset means. Bold and underline rank
methods only within the image+text
or image-only panel. Cross-panel differences are descriptive because the
public image+text methods receive target-derived synthetic captions unavailable
to image-only methods. Complete mean$\pm$sample standard deviation Dice and IoU
results are in Supplementary Table S14. The five \methodname{} entries use the
models from the ordered configuration analysis.\par
\end{minipage}
\end{table}
\fi

\ifnum\BmcTableSelection=7
\begin{table}[t]
\centering
\caption{Equal-dataset macro segmentation accuracy on the five public
benchmarks. Values are the arithmetic mean of the three-seed dataset means in
Table~\ref{tab:bspc_v3_models}. Bold and underline rank
models within modality. Cross-modality differences are descriptive because the
image+text models receive target-derived synthetic captions unavailable to the
image-only models.}
\label{tab:public5_macro_accuracy}
\small
\setlength{\tabcolsep}{6pt}
\begin{tabular}{lcc}
\toprule
Method & Dice (\%) & IoU (\%) \\
\midrule
\multicolumn{3}{l}{\textit{Image+text methods}} \\
\textbf{\methodname{} (Ours)} & $\mathbf{87.87}$ & $\mathbf{80.97}$ \\
CausalCLIPSeg Baseline & $\underline{86.53}$ & $\underline{79.32}$ \\
MedCLIPSeg & $86.11$ & $78.27$ \\
LViT & $84.85$ & $77.37$ \\
\midrule
\multicolumn{3}{l}{\textit{Image-only contextual references}} \\
nnU-Net & $\mathbf{84.96}$ & $\mathbf{78.23}$ \\
U-Mamba & $\underline{84.73}$ & $\underline{77.94}$ \\
U-KAN & $84.38$ & $77.02$ \\
U-Net++ & $84.10$ & $76.65$ \\
RollingUNet & $82.40$ & $74.62$ \\
U-Net & $79.62$ & $71.20$ \\
\bottomrule
\end{tabular}
\end{table}
\fi

\ifnum\BmcTableSelection=6
\begin{table}[htbp]
\centering
\caption{Exploratory three-seed continuous-location recomputation analysis.
\locrecompute{}
recomputes coarse position phrases from the transformed training mask when a
mild affine, translation, or zoom is sampled. Values are mean$\pm$sample
standard deviation. Delta and the reported bootstrap interval are
\locrecompute{} minus its reference model in percentage points. Both runs use
the same transformation distribution but independently initialized augmentation
sequences, so the realized per-case transformations are not paired. These
rows are not pooled with the main ordered configuration analysis. $p_W$ is the exploratory, unadjusted two-sided
paired Wilcoxon $p$-value.}
\label{tab:continuous_location_diagnostic}
\small
\setlength{\tabcolsep}{4.2pt}
\begin{tabular}{@{}llrrrr@{}}
\toprule
Dataset & Metric & Reference model & Location Recompute & $\Delta$ pp [95\% CI] & $p_W$ \\
\midrule
BUSI & Dice & $85.74\pm0.80$ & $84.62\pm0.20$ & $-1.12$ [$-2.62,0.07$] & .004 \\
     & IoU  & $79.46\pm0.83$ & $78.06\pm0.64$ & $-1.40$ [$-2.70,-0.30$] & .002 \\
ClinicDB & Dice & $93.18\pm0.54$ & $93.06\pm0.72$ & $-0.13$ [$-1.25,1.27$] & .138 \\
         & IoU  & $87.77\pm0.82$ & $87.54\pm1.07$ & $-0.23$ [$-1.94,1.78$] & .117 \\
\bottomrule
\end{tabular}
\end{table}
\fi


Across three seeds, exact-plan rewriting changes Public-5 macro Dice by
$-0.09$ pp (95\% dataset-stratified bootstrap CI $-0.30$ to $0.10$ pp) and macro IoU
by $-0.13$ pp (95\% CI $-0.39$ to $0.10$ pp). Dataset-level Dice differences
range from $+0.19$ pp on BUSI to $-0.30$ pp on ClinicDB. BRISC and COVID-19 have
Holm-adjusted Wilcoxon support for small negative rank shifts; the paired-case
bootstrap interval excludes zero for BRISC but not COVID-19. The macro interval
spans zero. Taken together, the exact control provides no evidence that prompt
rewriting improves average segmentation accuracy and identifies a small adverse
BRISC shift under the executed policy.

Trigger frequency helps interpret these differences but does not explain them
alone. COVID-19 has the highest realized prompt-change rate
($45.75\pm0.06\%$) yet its rewrite effect is slightly negative. BUSI has a lower
change rate ($20.27\pm0.10\%$) and a small positive point estimate. Trigger
frequency alone therefore does not predict the accuracy contrast; prompt
specificity, visual semantics, and optimization interactions remain plausible
moderators.

\def\BmcTableSelection{3}
% Tables for the BMC Medical Imaging manuscript.
% Ordered by the evidence hierarchy: dataset audit, ordered configurations,
% semantic controls, contextual model comparison, and exploratory diagnostics.

\ifnum\BmcTableSelection=1
\begin{table}[htbp]
\centering
\caption{Image-level dataset partitions and realized prompt-change frequencies.}
\label{tab:bspc_dataset_audit}
\scriptsize
\setlength{\tabcolsep}{2.0pt}
\begin{tabular}{@{}lrrrlrrr@{}}
\toprule
Dataset & Train & Val. & Test & Prompt & Unique & Spatial (\%) & Change (\%) \\
\midrule
BUSI     & 624  & 78  & 78   & Target-derived & 198 & 82.85 & $20.27\pm0.10$ \\
ClinicDB & 490  & 61  & 61   & Target-derived & 202 & 100.00 & $28.48\pm0.21$ \\
BUS-BRA  & 1311 & 282 & 282  & Target-derived & 75  & 100.00 & $12.44\pm0.12$ \\
BRISC    & 3881 & 974 & 892  & Target-derived & 359 & 78.23 & $5.75\pm0.04$ \\
COVID-19 & 5716 & 1429& 2113 & Target-derived & 269 & 100.00 & $45.75\pm0.06$ \\
\bottomrule
\end{tabular}
\vspace{2pt}
\begin{minipage}{0.98\linewidth}
\footnotesize\raggedright
\textit{Notes:} Unique and spatial coverage use training prompts. Realized
change enumerates all case--epoch events under the fixed 100-epoch shared
augmentation policy and reports mean$\pm$sample standard deviation across seeds
123, 456, and 789. Public prompts are target-derived synthetic captions, not
clinical text.\par
\end{minipage}
\end{table}
\fi

\ifnum\BmcTableSelection=3
\begin{table}[htbp]
\centering
\caption{Inference-time prompt controls for trained \methodname{} models on
five public datasets.}
\label{tab:bspc_semantic_controls}
\scriptsize
\setlength{\tabcolsep}{4.5pt}
\begin{tabular}{@{}lrrrrr@{}}
\toprule
Dataset & True Dice & Shuffled Dice & Fixed Dice & Empty Dice & $\Delta$ pp [95\% CI] \\
\midrule
BUSI     & $85.78\pm0.45$ & $46.73\pm1.41$ & $16.69\pm0.02$ & $13.22\pm5.35$ & $+39.05$ [$29.85,48.49$] \\
ClinicDB & $92.07\pm0.71$ & $90.20\pm0.94$ & $88.31\pm1.31$ & $78.71\pm7.93$ & $+1.88$ [$0.07,4.22$] \\
BUS-BRA  & $90.24\pm0.14$ & $89.47\pm0.23$ & $88.87\pm0.28$ & $68.22\pm8.44$ & $+0.77$ [$0.45,1.15$] \\
BRISC    & $87.18\pm0.14$ & $70.89\pm2.12$ & $75.31\pm1.97$ & $58.89\pm5.30$ & $+16.29$ [$14.26,18.43$] \\
COVID-19 & $84.08\pm0.13$ & $68.80\pm0.82$ & $68.54\pm0.16$ & $55.81\pm3.82$ & $+15.29$ [$14.49,16.10$] \\
\bottomrule
\end{tabular}
\vspace{2pt}
\begin{minipage}{0.98\linewidth}
\footnotesize\raggedright
\textit{Notes:} No model is retrained; only the test prompt changes. Dice is
mean$\pm$sample standard deviation over three seeds. The last column is true
minus shuffled Dice with a paired-case bootstrap 95\% interval after
three-seed averaging. Holm-adjusted Wilcoxon tests and the bootstrap intervals
are positive on all five datasets.\par
\end{minipage}
\end{table}
\fi

\ifnum\BmcTableSelection=4
% Auto-generated by generate_medequiseg_factorial_manuscript_assets.py.
% Do not edit numerical cells manually.
\begin{table}[htbp]
\centering
\caption{Five-dataset macro-average accuracy across five retained ordered configurations.}
\label{tab:medequiseg_factorial}
\scriptsize
\setlength{\tabcolsep}{2.4pt}
\begin{tabular}{@{}lccccrr@{}}
\toprule
Configuration & Text enc. & ATConv & Aug. & Rewrite & Dice & IoU \\
\midrule
Base model & -- & -- & -- & -- & $86.41\pm0.35$ & $79.22\pm0.34$ \\
$+$ BioMedCLIP & \checkmark & -- & -- & -- & $86.67\pm0.08$ & $79.52\pm0.06$ \\
$+$ ATConv & -- & \checkmark & -- & -- & $86.63\pm0.32$ & $79.42\pm0.31$ \\
$+$ BioMedCLIP $+$ ATConv & \checkmark & \checkmark & -- & -- & \underline{$87.03\pm0.15$} & \underline{$79.85\pm0.23$} \\
\textbf{\methodname{} (Ours)} & \checkmark & \checkmark & \checkmark & \checkmark & \textbf{$87.87\pm0.17$} & \textbf{$80.97\pm0.13$} \\
\bottomrule
\end{tabular}
\vspace{2pt}
\begin{minipage}{0.98\linewidth}
\footnotesize\raggedright
\textit{Notes:} Five retained configurations, five datasets, and three seeds produce 75
training runs. Values are five-dataset macro-average percentages
(mean$\pm$sample standard deviation). Text enc. denotes pooled BioMedCLIP
conditioning; ATConv denotes the two projection-head operators active in the evaluated forward pass;
Aug. denotes the common shared augmentation policy. Shared augmentation and
prompt rewriting appear together only in the final configuration, so these are conditional configuration contrasts,
not independent component effects or complete interactions. Bold and underline
mark the best and second-best configurations. The complete per-dataset matrix
is in Supplementary Table S9.\par
\end{minipage}
\end{table}

\fi

\ifnum\BmcTableSelection=5
\begin{table}[htbp]
\centering
\caption{Per-dataset and equal-dataset macro Dice on five public benchmarks.}
\label{tab:bspc_v3_models}
\scriptsize
\setlength{\tabcolsep}{2.0pt}
\begin{tabular}{@{}lrrrrrr@{}}
\toprule
Method & BUSI & ClinicDB & BUS-BRA & BRISC & COVID-19 & Macro avg. \\
\midrule
\multicolumn{7}{l}{\textit{Image+text methods}} \\
MedCLIPSeg
 & $\mathbf{86.13}$ & $\underline{91.96}$ & $88.47$
 & $81.36$ & $\underline{82.62}$ & $86.11$ \\
CausalCLIPSeg
 & $84.62$ & $91.19$ & $\underline{88.87}$
 & $\underline{86.11}$ & $81.85$ & $\underline{86.53}$ \\
LViT
 & $81.81$ & $91.03$ & $87.66$
 & $85.24$ & $78.50$ & $84.85$ \\
\textbf{\methodname{} (Ours)}
 & $\underline{85.78}$ & $\mathbf{92.07}$ & $\mathbf{90.24}$
 & $\mathbf{87.18}$ & $\mathbf{84.08}$ & $\mathbf{87.87}$ \\
\midrule
\multicolumn{7}{l}{\textit{Image-only contextual references}} \\
RollingUNet
 & $72.36$ & $87.31$ & $88.00$
 & $85.03$ & $79.32$ & $82.40$ \\
U-Net
 & $61.31$ & $85.61$ & $87.35$
 & $85.00$ & $78.81$ & $79.62$ \\
U-Net++
 & $76.21$ & $89.77$ & $88.59$
 & $\underline{85.94}$ & $\mathbf{80.00}$ & $84.10$ \\
nnU-Net
 & $74.99$ & $\mathbf{93.62}$ & $\underline{89.80}$
 & $\mathbf{87.30}$ & $79.07$ & $\mathbf{84.96}$ \\
U-Mamba
 & $\underline{76.63}$ & $\underline{92.21}$ & $\mathbf{90.34}$
 & $85.37$ & $79.08$ & $\underline{84.73}$ \\
U-KAN
 & $\mathbf{77.10}$ & $91.40$ & $88.56$
 & $85.45$ & $\underline{79.38}$ & $84.38$ \\
\bottomrule
\end{tabular}
\vspace{2pt}
\begin{minipage}{0.98\linewidth}
\footnotesize\raggedright
\textit{Notes:} Dataset values are three-seed mean Dice percentages; the macro
average is the arithmetic mean of the five dataset means. Bold and underline rank
methods only within the image+text
or image-only panel. Cross-panel differences are descriptive because the
public image+text methods receive target-derived synthetic captions unavailable
to image-only methods. Complete mean$\pm$sample standard deviation Dice and IoU
results are in Supplementary Table S14. The five \methodname{} entries use the
models from the ordered configuration analysis.\par
\end{minipage}
\end{table}
\fi

\ifnum\BmcTableSelection=7
\begin{table}[t]
\centering
\caption{Equal-dataset macro segmentation accuracy on the five public
benchmarks. Values are the arithmetic mean of the three-seed dataset means in
Table~\ref{tab:bspc_v3_models}. Bold and underline rank
models within modality. Cross-modality differences are descriptive because the
image+text models receive target-derived synthetic captions unavailable to the
image-only models.}
\label{tab:public5_macro_accuracy}
\small
\setlength{\tabcolsep}{6pt}
\begin{tabular}{lcc}
\toprule
Method & Dice (\%) & IoU (\%) \\
\midrule
\multicolumn{3}{l}{\textit{Image+text methods}} \\
\textbf{\methodname{} (Ours)} & $\mathbf{87.87}$ & $\mathbf{80.97}$ \\
CausalCLIPSeg Baseline & $\underline{86.53}$ & $\underline{79.32}$ \\
MedCLIPSeg & $86.11$ & $78.27$ \\
LViT & $84.85$ & $77.37$ \\
\midrule
\multicolumn{3}{l}{\textit{Image-only contextual references}} \\
nnU-Net & $\mathbf{84.96}$ & $\mathbf{78.23}$ \\
U-Mamba & $\underline{84.73}$ & $\underline{77.94}$ \\
U-KAN & $84.38$ & $77.02$ \\
U-Net++ & $84.10$ & $76.65$ \\
RollingUNet & $82.40$ & $74.62$ \\
U-Net & $79.62$ & $71.20$ \\
\bottomrule
\end{tabular}
\end{table}
\fi

\ifnum\BmcTableSelection=6
\begin{table}[htbp]
\centering
\caption{Exploratory three-seed continuous-location recomputation analysis.
\locrecompute{}
recomputes coarse position phrases from the transformed training mask when a
mild affine, translation, or zoom is sampled. Values are mean$\pm$sample
standard deviation. Delta and the reported bootstrap interval are
\locrecompute{} minus its reference model in percentage points. Both runs use
the same transformation distribution but independently initialized augmentation
sequences, so the realized per-case transformations are not paired. These
rows are not pooled with the main ordered configuration analysis. $p_W$ is the exploratory, unadjusted two-sided
paired Wilcoxon $p$-value.}
\label{tab:continuous_location_diagnostic}
\small
\setlength{\tabcolsep}{4.2pt}
\begin{tabular}{@{}llrrrr@{}}
\toprule
Dataset & Metric & Reference model & Location Recompute & $\Delta$ pp [95\% CI] & $p_W$ \\
\midrule
BUSI & Dice & $85.74\pm0.80$ & $84.62\pm0.20$ & $-1.12$ [$-2.62,0.07$] & .004 \\
     & IoU  & $79.46\pm0.83$ & $78.06\pm0.64$ & $-1.40$ [$-2.70,-0.30$] & .002 \\
ClinicDB & Dice & $93.18\pm0.54$ & $93.06\pm0.72$ & $-0.13$ [$-1.25,1.27$] & .138 \\
         & IoU  & $87.77\pm0.82$ & $87.54\pm1.07$ & $-0.23$ [$-1.94,1.78$] & .117 \\
\bottomrule
\end{tabular}
\end{table}
\fi


The semantic controls further separate prompt use from rewrite benefit. True
minus shuffled Dice ranges from $+0.77$ pp on BUS-BRA and $+1.88$ pp on
ClinicDB to $+39.05$ pp on BUSI. Holm-adjusted Wilcoxon tests support a
true--shuffled difference on all five public datasets, and the paired-case
bootstrap intervals are positive throughout; ClinicDB has the smallest interval
($+0.07$ to $+4.22$ pp). This demonstrates sensitivity to privileged case-matched caption
identity rather than clinical-report understanding or specifically spatial
grounding, because the public captions contain target-derived attributes.

Presence/class matching substantially attenuates but does not eliminate this
effect. True text exceeds the stratified derangement by $+5.63$ Dice pp in the
equal-dataset Public-5 macro (95\% CI $+4.71$ to $+6.67$), compared with
$+14.66$ pp under unstratified shuffling. Four of five dataset bootstrap
intervals are positive; Holm-adjusted rank support remains for BUSI, BRISC, and
COVID-19, whereas BUS-BRA is neutral. Restricting the analysis to the 3,413
cases whose text string actually changes gives $+5.87$ pp (95\% CI $+4.87$ to
$+7.03$). For the 474 recoverable COVID-19 subjects, equal-subject averaging
gives $+22.42$ pp with a subject-cluster interval of $+21.12$ to $+23.80$.
These post hoc results show that coarse class/presence mismatch explains a large
share of the original shuffled gap while leaving a heterogeneous residual
association with target-derived prompt identity (Supplementary Table S13).

\begin{figure}[!htbp]
\centering
\includegraphics[width=\linewidth]{protocol_v3_semantic_control_audit.pdf}
\caption{Fixed-checkpoint true-versus-shuffled prompt audit. One non-empty case
nearest the 75th percentile of the three-seed per-case Dice difference is shown
for each dataset that passes the Holm-adjusted rank screen. Predictions are
three-seed majority votes. The left-hand selection delta is the mean paired
per-case difference across seeds and therefore need not equal the difference
between the displayed majority-vote Dice values. Each row uses one shared context crop determined by
the union of the ground truth and both predictions. Green denotes true positive,
red false positive, and blue false negative. Selection was deterministic; exact
prompts, case identifiers, and uncropped paths are archived in the selection CSV.}
\label{fig:semantic_control_audit}
\end{figure}

\FloatBarrier
\subsection{Per-dataset comparisons and qualitative error analysis}

\def\BmcTableSelection{5}
% Tables for the BMC Medical Imaging manuscript.
% Ordered by the evidence hierarchy: dataset audit, ordered configurations,
% semantic controls, contextual model comparison, and exploratory diagnostics.

\ifnum\BmcTableSelection=1
\begin{table}[htbp]
\centering
\caption{Image-level dataset partitions and realized prompt-change frequencies.}
\label{tab:bspc_dataset_audit}
\scriptsize
\setlength{\tabcolsep}{2.0pt}
\begin{tabular}{@{}lrrrlrrr@{}}
\toprule
Dataset & Train & Val. & Test & Prompt & Unique & Spatial (\%) & Change (\%) \\
\midrule
BUSI     & 624  & 78  & 78   & Target-derived & 198 & 82.85 & $20.27\pm0.10$ \\
ClinicDB & 490  & 61  & 61   & Target-derived & 202 & 100.00 & $28.48\pm0.21$ \\
BUS-BRA  & 1311 & 282 & 282  & Target-derived & 75  & 100.00 & $12.44\pm0.12$ \\
BRISC    & 3881 & 974 & 892  & Target-derived & 359 & 78.23 & $5.75\pm0.04$ \\
COVID-19 & 5716 & 1429& 2113 & Target-derived & 269 & 100.00 & $45.75\pm0.06$ \\
\bottomrule
\end{tabular}
\vspace{2pt}
\begin{minipage}{0.98\linewidth}
\footnotesize\raggedright
\textit{Notes:} Unique and spatial coverage use training prompts. Realized
change enumerates all case--epoch events under the fixed 100-epoch shared
augmentation policy and reports mean$\pm$sample standard deviation across seeds
123, 456, and 789. Public prompts are target-derived synthetic captions, not
clinical text.\par
\end{minipage}
\end{table}
\fi

\ifnum\BmcTableSelection=3
\begin{table}[htbp]
\centering
\caption{Inference-time prompt controls for trained \methodname{} models on
five public datasets.}
\label{tab:bspc_semantic_controls}
\scriptsize
\setlength{\tabcolsep}{4.5pt}
\begin{tabular}{@{}lrrrrr@{}}
\toprule
Dataset & True Dice & Shuffled Dice & Fixed Dice & Empty Dice & $\Delta$ pp [95\% CI] \\
\midrule
BUSI     & $85.78\pm0.45$ & $46.73\pm1.41$ & $16.69\pm0.02$ & $13.22\pm5.35$ & $+39.05$ [$29.85,48.49$] \\
ClinicDB & $92.07\pm0.71$ & $90.20\pm0.94$ & $88.31\pm1.31$ & $78.71\pm7.93$ & $+1.88$ [$0.07,4.22$] \\
BUS-BRA  & $90.24\pm0.14$ & $89.47\pm0.23$ & $88.87\pm0.28$ & $68.22\pm8.44$ & $+0.77$ [$0.45,1.15$] \\
BRISC    & $87.18\pm0.14$ & $70.89\pm2.12$ & $75.31\pm1.97$ & $58.89\pm5.30$ & $+16.29$ [$14.26,18.43$] \\
COVID-19 & $84.08\pm0.13$ & $68.80\pm0.82$ & $68.54\pm0.16$ & $55.81\pm3.82$ & $+15.29$ [$14.49,16.10$] \\
\bottomrule
\end{tabular}
\vspace{2pt}
\begin{minipage}{0.98\linewidth}
\footnotesize\raggedright
\textit{Notes:} No model is retrained; only the test prompt changes. Dice is
mean$\pm$sample standard deviation over three seeds. The last column is true
minus shuffled Dice with a paired-case bootstrap 95\% interval after
three-seed averaging. Holm-adjusted Wilcoxon tests and the bootstrap intervals
are positive on all five datasets.\par
\end{minipage}
\end{table}
\fi

\ifnum\BmcTableSelection=4
% Auto-generated by generate_medequiseg_factorial_manuscript_assets.py.
% Do not edit numerical cells manually.
\begin{table}[htbp]
\centering
\caption{Five-dataset macro-average accuracy across five retained ordered configurations.}
\label{tab:medequiseg_factorial}
\scriptsize
\setlength{\tabcolsep}{2.4pt}
\begin{tabular}{@{}lccccrr@{}}
\toprule
Configuration & Text enc. & ATConv & Aug. & Rewrite & Dice & IoU \\
\midrule
Base model & -- & -- & -- & -- & $86.41\pm0.35$ & $79.22\pm0.34$ \\
$+$ BioMedCLIP & \checkmark & -- & -- & -- & $86.67\pm0.08$ & $79.52\pm0.06$ \\
$+$ ATConv & -- & \checkmark & -- & -- & $86.63\pm0.32$ & $79.42\pm0.31$ \\
$+$ BioMedCLIP $+$ ATConv & \checkmark & \checkmark & -- & -- & \underline{$87.03\pm0.15$} & \underline{$79.85\pm0.23$} \\
\textbf{\methodname{} (Ours)} & \checkmark & \checkmark & \checkmark & \checkmark & \textbf{$87.87\pm0.17$} & \textbf{$80.97\pm0.13$} \\
\bottomrule
\end{tabular}
\vspace{2pt}
\begin{minipage}{0.98\linewidth}
\footnotesize\raggedright
\textit{Notes:} Five retained configurations, five datasets, and three seeds produce 75
training runs. Values are five-dataset macro-average percentages
(mean$\pm$sample standard deviation). Text enc. denotes pooled BioMedCLIP
conditioning; ATConv denotes the two projection-head operators active in the evaluated forward pass;
Aug. denotes the common shared augmentation policy. Shared augmentation and
prompt rewriting appear together only in the final configuration, so these are conditional configuration contrasts,
not independent component effects or complete interactions. Bold and underline
mark the best and second-best configurations. The complete per-dataset matrix
is in Supplementary Table S9.\par
\end{minipage}
\end{table}

\fi

\ifnum\BmcTableSelection=5
\begin{table}[htbp]
\centering
\caption{Per-dataset and equal-dataset macro Dice on five public benchmarks.}
\label{tab:bspc_v3_models}
\scriptsize
\setlength{\tabcolsep}{2.0pt}
\begin{tabular}{@{}lrrrrrr@{}}
\toprule
Method & BUSI & ClinicDB & BUS-BRA & BRISC & COVID-19 & Macro avg. \\
\midrule
\multicolumn{7}{l}{\textit{Image+text methods}} \\
MedCLIPSeg
 & $\mathbf{86.13}$ & $\underline{91.96}$ & $88.47$
 & $81.36$ & $\underline{82.62}$ & $86.11$ \\
CausalCLIPSeg
 & $84.62$ & $91.19$ & $\underline{88.87}$
 & $\underline{86.11}$ & $81.85$ & $\underline{86.53}$ \\
LViT
 & $81.81$ & $91.03$ & $87.66$
 & $85.24$ & $78.50$ & $84.85$ \\
\textbf{\methodname{} (Ours)}
 & $\underline{85.78}$ & $\mathbf{92.07}$ & $\mathbf{90.24}$
 & $\mathbf{87.18}$ & $\mathbf{84.08}$ & $\mathbf{87.87}$ \\
\midrule
\multicolumn{7}{l}{\textit{Image-only contextual references}} \\
RollingUNet
 & $72.36$ & $87.31$ & $88.00$
 & $85.03$ & $79.32$ & $82.40$ \\
U-Net
 & $61.31$ & $85.61$ & $87.35$
 & $85.00$ & $78.81$ & $79.62$ \\
U-Net++
 & $76.21$ & $89.77$ & $88.59$
 & $\underline{85.94}$ & $\mathbf{80.00}$ & $84.10$ \\
nnU-Net
 & $74.99$ & $\mathbf{93.62}$ & $\underline{89.80}$
 & $\mathbf{87.30}$ & $79.07$ & $\mathbf{84.96}$ \\
U-Mamba
 & $\underline{76.63}$ & $\underline{92.21}$ & $\mathbf{90.34}$
 & $85.37$ & $79.08$ & $\underline{84.73}$ \\
U-KAN
 & $\mathbf{77.10}$ & $91.40$ & $88.56$
 & $85.45$ & $\underline{79.38}$ & $84.38$ \\
\bottomrule
\end{tabular}
\vspace{2pt}
\begin{minipage}{0.98\linewidth}
\footnotesize\raggedright
\textit{Notes:} Dataset values are three-seed mean Dice percentages; the macro
average is the arithmetic mean of the five dataset means. Bold and underline rank
methods only within the image+text
or image-only panel. Cross-panel differences are descriptive because the
public image+text methods receive target-derived synthetic captions unavailable
to image-only methods. Complete mean$\pm$sample standard deviation Dice and IoU
results are in Supplementary Table S14. The five \methodname{} entries use the
models from the ordered configuration analysis.\par
\end{minipage}
\end{table}
\fi

\ifnum\BmcTableSelection=7
\begin{table}[t]
\centering
\caption{Equal-dataset macro segmentation accuracy on the five public
benchmarks. Values are the arithmetic mean of the three-seed dataset means in
Table~\ref{tab:bspc_v3_models}. Bold and underline rank
models within modality. Cross-modality differences are descriptive because the
image+text models receive target-derived synthetic captions unavailable to the
image-only models.}
\label{tab:public5_macro_accuracy}
\small
\setlength{\tabcolsep}{6pt}
\begin{tabular}{lcc}
\toprule
Method & Dice (\%) & IoU (\%) \\
\midrule
\multicolumn{3}{l}{\textit{Image+text methods}} \\
\textbf{\methodname{} (Ours)} & $\mathbf{87.87}$ & $\mathbf{80.97}$ \\
CausalCLIPSeg Baseline & $\underline{86.53}$ & $\underline{79.32}$ \\
MedCLIPSeg & $86.11$ & $78.27$ \\
LViT & $84.85$ & $77.37$ \\
\midrule
\multicolumn{3}{l}{\textit{Image-only contextual references}} \\
nnU-Net & $\mathbf{84.96}$ & $\mathbf{78.23}$ \\
U-Mamba & $\underline{84.73}$ & $\underline{77.94}$ \\
U-KAN & $84.38$ & $77.02$ \\
U-Net++ & $84.10$ & $76.65$ \\
RollingUNet & $82.40$ & $74.62$ \\
U-Net & $79.62$ & $71.20$ \\
\bottomrule
\end{tabular}
\end{table}
\fi

\ifnum\BmcTableSelection=6
\begin{table}[htbp]
\centering
\caption{Exploratory three-seed continuous-location recomputation analysis.
\locrecompute{}
recomputes coarse position phrases from the transformed training mask when a
mild affine, translation, or zoom is sampled. Values are mean$\pm$sample
standard deviation. Delta and the reported bootstrap interval are
\locrecompute{} minus its reference model in percentage points. Both runs use
the same transformation distribution but independently initialized augmentation
sequences, so the realized per-case transformations are not paired. These
rows are not pooled with the main ordered configuration analysis. $p_W$ is the exploratory, unadjusted two-sided
paired Wilcoxon $p$-value.}
\label{tab:continuous_location_diagnostic}
\small
\setlength{\tabcolsep}{4.2pt}
\begin{tabular}{@{}llrrrr@{}}
\toprule
Dataset & Metric & Reference model & Location Recompute & $\Delta$ pp [95\% CI] & $p_W$ \\
\midrule
BUSI & Dice & $85.74\pm0.80$ & $84.62\pm0.20$ & $-1.12$ [$-2.62,0.07$] & .004 \\
     & IoU  & $79.46\pm0.83$ & $78.06\pm0.64$ & $-1.40$ [$-2.70,-0.30$] & .002 \\
ClinicDB & Dice & $93.18\pm0.54$ & $93.06\pm0.72$ & $-0.13$ [$-1.25,1.27$] & .138 \\
         & IoU  & $87.77\pm0.82$ & $87.54\pm1.07$ & $-0.23$ [$-1.94,1.78$] & .117 \\
\bottomrule
\end{tabular}
\end{table}
\fi


The broad model table is contextual. \methodname{} has a higher mean Dice than
reproduced MedCLIPSeg on ClinicDB ($+0.11$ pp), BUS-BRA ($+1.77$ pp), BRISC
($+5.82$ pp), and COVID-19 ($+1.46$ pp), and a lower mean on BUSI ($-0.35$ pp).
Relative to the strongest image-only mean selected within each dataset,
\methodname{} is higher on BUSI ($+8.68$ Dice pp) and COVID-19 ($+4.08$ pp), and
lower by $0.10$--$1.55$ pp on the other public datasets. These
cross-modality differences are descriptive, not evidence that privileged
target-derived text yields a fair deployment advantage.

The COVID-19 grouped train--validation sensitivity preserved the fixed test set
and removed all overlap among the 2,130 recoverable subject tokens. In this
separate sensitivity checkpoint set, mean Dice was $83.99\pm0.07\%$ for the
\methodname{} configuration, compared with $84.18\pm0.13\%$ for its
configuration-matched original-split rerun, a change of $-0.19$ pp. The
corresponding no-rewrite and U-Net++ changes were $+0.00$ and $+0.14$ pp. Within
the grouped-validation checkpoint set, rewriting is $-0.29$ Dice pp below
no-rewrite (paired-case 95\% CI $-0.55$ to $-0.03$); the equal-subject effect is
$-0.40$ pp with a subject-cluster interval of $-0.58$ to $-0.02$ and Wilcoxon
$p<.001$. Thus, the known recoverable-subject
overlap has little effect on the comparison with U-Net++, while the grouped
sensitivity reinforces rather than reverses the absence of a rewrite benefit. However,
rows without recoverable subject tokens cannot be certified as patient grouped
(Supplementary Table S10).

The mask-presence analysis (Supplementary Table S8) shows that BUSI contains 65
non-empty and 13 empty test masks. \methodname{} and U-KAN have non-empty Dice values of
$82.94\pm0.55\%$ and $79.70\pm0.48\%$, a difference of $+3.24$ pp, whereas
the overall difference is $+8.68$ pp. Empty specificity is $100.00\%$ for \methodname{}
and $64.10\pm17.76\%$ for U-KAN. Empty-mask cases therefore amplify, but do
not fully account for, the BUSI difference. BRISC has the only other mixed
test split (753 non-empty and 139 empty masks); the complete subgroup audit is
reported rather than using empty cases as implicit evidence of lesion quality.

The canonical boundary results in Supplementary Tables S1--S5 follow the same
dataset-dependent pattern. Because one-sided empty cases have infinite distance,
finite HD95 and ASSD are interpreted only with the failure rate.

% Auto-generated by paper/analysis/generate_complexity_table.py
\begin{table}[!htbp]
\centering
\caption{Instantiated model size and synchronized single-image inference cost.}
\label{tab:complexity}
\scriptsize
\setlength{\tabcolsep}{2.0pt}
\begin{tabular}{@{}lrrrrr@{}}
\toprule
Method & Params (M) & Trainable (M) & ms/image & Images/s & Peak MB \\
\midrule
U-Net++$^\dagger$ & 2.41 & 2.41 & 6.80$\pm$0.09 & 146.95 & 123 \\
RollingUNet$^\dagger$ & 1.78 & 1.78 & 26.40$\pm$0.65 & 37.88 & 31 \\
LViT & 39.93 & 39.93 & 11.35$\pm$0.58 & 88.08 & 278 \\
CausalCLIPSeg Baseline & 165.26 & 165.26 & 18.83$\pm$0.53 & 53.10 & 895 \\
MedCLIPSeg ($S=1$) & 214.58 & 18.68 & 55.44$\pm$6.97 & 18.04 & 1814 \\
MedCLIPSeg ($S=30$) & 214.58 & 18.68 & 816.41$\pm$12.91 & 1.22 & 1825 \\
\textbf{\methodname{} (Ours)} & 159.78 & 159.78 & 15.03$\pm$0.82 & 66.55 & 874 \\
\bottomrule
\end{tabular}
\vspace{2pt}
\begin{minipage}{0.98\linewidth}
\footnotesize\raggedright
\textit{Notes:} Parameter columns count all instantiated modules, not only
output-connected subgraphs; trainable counts follow the saved model's parameter
training-state flags. Timing uses batch size 1, FP32, $224\times224$ input,
and an NVIDIA GeForce RTX 4090 and is reported as mean$\pm$standard deviation
after warm-up; data loading and post-processing are excluded. Cached text
lookup is included for \methodname{} and LViT, whereas MedCLIPSeg runs its
online UniMedCLIP/BiomedBERT towers. $S$ is the requested number of
probabilistic samples. A dagger marks image-only methods.\par
\end{minipage}
\end{table}


\methodname{} instantiates 159.78 million parameters and requires
$15.03\pm0.82$ ms per image. The training-time policy adds no inference
parameters, and \norewrite{} has the same inference graph. These measurements describe
implementation cost rather than an architecture-controlled speed comparison.

\begin{figure}[!htbp]
\centering
\includegraphics[width=\linewidth]{protocol_v3_qualitative_audit.pdf}
\caption{Frozen-rule qualitative audit. BUSI and COVID-19 cases are nearest the
pre-specified 90th percentile of the three-seed \methodname{} minus image-only Dice
difference; ClinicDB and BUS-BRA cases are nearest the 10th percentile.
Predictions are pixel-wise majority votes. The left-hand selection delta is the
mean paired per-case difference across seeds and therefore need not equal the
difference between the displayed majority-vote Dice values. Green is true positive, red false
positive, and blue false negative. Within each row, all columns use one context
crop determined by the union of the ground truth and both predictions. The
deterministic case and prompt record is provided in the results archive;
examples are not additional tests.}
\label{fig:qualitative_audit}
\end{figure}

The pre-specified upper- and lower-tail examples in
Figure~\ref{fig:qualitative_audit} make the claim boundary visible: selected BUSI
and COVID-19 cases favor \methodname{}, whereas the ClinicDB and BUS-BRA cases
expose false-positive or false-negative failure modes.

\FloatBarrier
\section{Discussion}
The principal empirical result is higher accuracy for the complete configuration
than for the base multimodal pathway: $+1.46$ Dice pp (95\% CI $+0.95$ to
$+1.99$) and $+1.74$ IoU pp (95\% CI $+1.13$ to $+2.34$) in the Public-5
ordered ablation. Because the contrast was selected retrospectively and the
shared policy was a follow-up amendment on the fixed V3 manifests and evaluator,
it supports an association on this benchmark panel rather than a prospectively
confirmed universal effect. \methodname{} also has the highest equal-dataset
macro Dice and IoU among the reproduced models, but this does not mean that it
is the best method on every dataset.
The strongest image-only result is higher on ClinicDB, BUS-BRA, and BRISC,
whereas \methodname{} is higher on BUSI and COVID-19. Reporting both macro and
per-dataset results is therefore necessary: the macro captures consistency
across the benchmark panel, while the dataset table exposes modality-specific
failure and competitive margins.

The component study localizes this gain. BioMedCLIP and ATConv individually
change the Public-5 macro by only $+0.26$ and $+0.22$ Dice pp relative to the
base pathway; using both gives $+0.62$ pp. The shared plan without rewriting
reaches $+1.55$ pp, while the complete configuration reaches $+1.46$ pp. The
post hoc Base-plus-plan control independently reaches $+1.42$ pp, while the
complete-model residual over that control is $+0.04$ pp (95\% CI $-0.40$ to
$+0.41$). The shared policy therefore explains nearly all of the observed macro
point improvement; the data do not establish an additional aggregate accuracy
benefit from the complete BioMedCLIP--ATConv--rewrite stack beyond that policy.
The runtime audit confirms that ATConv was not globally disabled: two projector
operators execute and produce non-zero outputs. It also confirms that the other
two registered candidates are dormant, so these results support only the
effective two-projector adapter and not an ATConv4 claim.

Prompt reliability remains an interpretive control rather than the primary
performance claim. A shared plan guarantees image--mask alignment; the language
map addresses prompt validity under a subset of spatial transforms; and semantic
controls test whether the checkpoint depends on case-matched text. The exact-plan
control is essential because comparison with an unaugmented baseline cannot
distinguish language consistency from ordinary regularization. Likewise,
segmentation accuracy without shuffled or empty prompts cannot establish reliable
text use.

The semantic-control pattern constrains interpretation. Large true--shuffled
gaps on all five public sets show that prompt identity affects the prediction.
After matching target presence and diagnostic class, the Public-5 gap falls from
$+14.66$ to $+5.63$ Dice pp. This attenuation identifies coarse class/presence
confounding in the unstratified control, while the remaining effect is
concentrated in BUSI and COVID-19 and is nearly absent on BUS-BRA. The residual
therefore supports sensitivity to privileged case-specific prompt content, not
a general claim of language grounding.
The constant-text training control also differs from the complete model by
$+1.90$ Dice pp in the Public-5 macro, but only BUSI and COVID-19 have
dataset-level intervals excluding zero. Target-derived captions therefore make
this a concentrated privileged synthetic referring effect rather than
clinical-report understanding.

The result is consistent with the broader paired-augmentation literature rather
than a priority claim. Augment the Pairs rewrites predefined caption keywords
during horizontal flips, while RBTCA builds a region-grounded modality-aware
representation before augmentation \cite{augmentpairs,rbtca}. Our added evidence
is chiefly diagnostic: a deterministic exact-plan no-rewrite comparison shows
that semantic correction need not improve overlap accuracy even when the shared
augmentation policy does.

The prompt-design evidence from DescriptorMedSAM and the explicit spatial
reasoning in SSA reinforce that location language can be useful when it is both
informative and correctly grounded \cite{descriptormedsam,ssa}. NTP-MRISeg
further shows that text can be integrated through a unified generative objective
rather than a conditional decoder \cite{ntpmriseg}. Their published values are
not inserted into our quantitative tables because their datasets, prompt
construction, splits, architectures, and evaluation recipes are not identical to
Protocol V3. Accordingly, they contextualize the design space but do not serve as
same-protocol evidence for or against \methodname{}.

The BRISC result requires particular care because duplicate sanitation changes
the effective sample count. The sanitation rule is fixed before training and
applied to all compared methods.

\subsection{Limitations and next evidence}
First, the V3 hash identifies fixed data and evaluation choices, not the complete
training policy. The richer R11 shared plan changes the original dataset routing
and adds affine, translation, and zoom, so the 90-cell matrix is a frozen
follow-up ablation rather than a strict execution of the original July 10
augmentation specification. Existing run-configuration and code hashes preserve
the identity of the executed batch, but the recorded Git commit omits several
orchestration/evaluation files and the currently available release copies do not
reproduce the aggregate execution-code digest. Exact clean-clone reconstruction
of the 90-cell source snapshot is therefore not demonstrated. The original
completion digest also omitted command-line fields such as batch size and worker
count; archived command logs are needed to recover them. Second, four ATConv
modules are present in checkpoint metadata,
but only the two projector modules enter the released forward graph. The reported
parameter count includes the two dormant neck modules, and no ATConv4 claim is
made. Third, the exact three-seed rewrite control shows small heterogeneous
effects and a macro bootstrap interval spanning zero. BRISC has a small negative
paired-case interval, and the grouped COVID-19 sensitivity is also negative;
neither analysis supports an accuracy-benefit claim for rewriting. Fourth, the
current framework audits reliability at dataset and checkpoint level but does
not contain a learned per-case confidence gate or automatic image-only fallback.
Fifth, continuous translation and zoom in the reported \methodname{} results do not recompute coarse
location labels. The completed three-seed \locrecompute{} diagnostic decreases
BUSI performance and is neutral on ClinicDB, but it is mask derived and does not
share exact realized transform draws with the main configuration; it cannot
establish the effect of a clinically supplied continuous-language update. Sixth,
the public prompts contain attributes extracted from the target masks or labels.
This controlled referring benchmark is useful for testing equivariance, but it
gives image+text methods privileged target metadata and cannot establish clinical
utility or a fair deployment advantage over image-only models. The post hoc
stratified derangement controls only coarse presence and diagnostic
class derived from released identifiers/prompts; it does not match morphology,
location, wording template, or other target-derived attributes and therefore
cannot isolate clinical language grounding. None of the prompts were authored
prospectively by clinicians. Seventh, the
component study is complete but is an ordered ablation rather than a
full crossing of every augmentation-policy interaction. Eighth, although six image-only architectures are evaluated with
complete three-seed coverage, architectural recipes are not forced to share input
size, loss, or augmentation; the comparison is controlled at the data and
evaluation levels rather than as an identical-training ablation. The broad model
table is therefore contextual, with rankings separated by modality. Ninth, the
five public datasets differ in modality, sample count, prompt diversity, and
spatial-language prevalence, so an aggregate average would obscure meaningful
heterogeneity. Tenth, the public benchmark manifests do not generally expose
patient or acquisition-sequence identifiers. The fixed test folders are
image-disjoint, and the recoverable COVID-19 subject IDs do not cross into test,
but patient-level independence cannot be established for every dataset from the
released metadata. For COVID-19, the grouped train--validation sensitivity
removes all overlap among 2,130 recoverable subject tokens and changes
\methodname{} Dice by only $-0.19$ pp, but 2,951 rows without recoverable tokens
remain case-level groups. The reported-checkpoint subject-cluster sensitivity
preserves the direction of the large U-Net++ and shuffled-prompt contrasts,
while the small rewrite contrast remains uncertain; comparable
clustered inference is not possible for the other public datasets from the
released metadata. ClinicDB sequence independence also cannot be verified because
sequence identifiers are absent.

The current manuscript therefore limits its primary claim to a fixed-data-and-evaluation
benchmark association led by the shared training policy. Prospective clinician-authored text and independent
cross-domain testing remain future validation rather than prerequisites for the
reported benchmark comparison.
Complete boundary-metric tables, model complexity, and qualitative failure
panels are provided in the supplement and main manuscript.

\section{Conclusions}
This study provides a 90-cell, three-seed audit of shared augmentation and
privileged prompts in multimodal medical image segmentation. The complete
configuration achieves Public-5 macro Dice of $87.87\%$ and IoU of $80.97\%$ and
differs from the horizontal-flip Base reference by $+1.46$ Dice pp (95\% CI $+0.95$ to
$+1.99$) and $+1.74$ IoU pp (95\% CI $+1.13$ to $+2.34$). The shared-plan
no-rewrite row is numerically best, and the complete configuration differs from
the post hoc Base-plus-plan control by only $+0.04$ Dice pp (95\% CI $-0.40$ to
$+0.41$). The supported accuracy claim is therefore attached primarily to the
shared follow-up training policy, not to rewriting or to an ATConv4 mechanism.
Case-specific target-derived text affects aggregate predictions after coarse
presence/class matching, but the effect is markedly attenuated and does not
establish clinical-report understanding. The protocol boundary, effective
two-projector ATConv graph, negative rewrite evidence, and privileged-prompt
control are part of the result, not implementation details to be hidden from
interpretation.
